%Version 3.1 December 2024
% See section 11 of the User Manual for version history
%
%%%%%%%%%%%%%%%%%%%%%%%%%%%%%%%%%%%%%%%%%%%%%%%%%%%%%%%%%%%%%%%%%%%%%%
%%                                                                 %%
%% Please do not use \input{...} to include other tex files.       %%
%% Submit your LaTeX manuscript as one .tex document.              %%
%%                                                                 %%
%% All additional figures and files should be attached             %%
%% separately and not embedded in the \TeX\ document itself.       %%
%%                                                                 %%
%%%%%%%%%%%%%%%%%%%%%%%%%%%%%%%%%%%%%%%%%%%%%%%%%%%%%%%%%%%%%%%%%%%%%%

%% ---------------------------------------------------------------------------------------
%% DRAFT STATUS (Track A). Sections 2-5, 9.1-9.2 and Appendices A-E are drafted from the
%% committed Phase 1 results. Sections marked TRACK B wait for the multilingual user study.
%% Every number traces to PHASE1_RESULTS.md, PHASE1_RESULTS_RAG.md, PHASE1_COMPARISON.md,
%% PHASE1_SCORECARD.md or PHASE1_PREREGISTRATION.md on branch phase1-test-bed.
%% LANGUAGE_TEST_RESULTS.md (13/7/2 tiers, one lesson per subject) is superseded: do not
%% quote it except as the pilot mentioned in Section 4.
%% Long tables (22-language matrices, lesson bank, language codes) were emitted by a script
%% from those files and lessons.json, not typed by hand.
%% ---------------------------------------------------------------------------------------

\documentclass[pdflatex,sn-mathphys-num]{sn-jnl}% Math and Physical Sciences Numbered Reference Style

\usepackage{graphicx}%
\usepackage{multirow}%
\usepackage{amsmath,amssymb,amsfonts}%
\usepackage{amsthm}%
\usepackage{mathrsfs}%
\usepackage[title]{appendix}%
\usepackage{xcolor}%
\usepackage{textcomp}%
\usepackage{manyfoot}%
\usepackage{booktabs}%
\usepackage{algorithm}%
\usepackage{algorithmicx}%
\usepackage{algpseudocode}%
\usepackage{listings}%
\usepackage{longtable}%
\usepackage{array}%

\newcommand{\todo}[1]{\textcolor{red}{[TODO: #1]}}
\newcommand{\trackb}[1]{\textcolor{blue}{[TRACK B: #1]}}
\lstset{basicstyle=\ttfamily\footnotesize, breaklines=true, columns=fullflexible, frame=single}

\theoremstyle{thmstyleone}%
\newtheorem{theorem}{Theorem}%
\newtheorem{proposition}[theorem]{Proposition}%
\theoremstyle{thmstyletwo}%
\newtheorem{example}{Example}%
\newtheorem{remark}{Remark}%
\theoremstyle{thmstylethree}%
\newtheorem{definition}{Definition}%

\raggedbottom

\begin{document}

\title[Teaching an AI Learner in Your Own Language]{Teaching an AI Learner in Your Own Language: A Translation Boundary for Learning-by-Teaching Tutors in 22 Indian Languages and a Comparison of Two AI Learners}

\author[1]{\fnm{First} \sur{Author}}\email{iauthor@gmail.com}
\author[1]{\fnm{Second} \sur{Author}}\email{iiauthor@gmail.com}
\equalcont{These authors contributed equally to this work.}
\author[1]{\fnm{Third} \sur{Author}}\email{iiiauthor@gmail.com}
\equalcont{These authors contributed equally to this work.}
\author[1]{\fnm{Dhruv} \sur{Kumar}}\email{dhruv.kumar@pilani.bits-pilani.ac.in}
\author*[1]{\fnm{Jagat Sesh} \sur{Challa}}\email{jagatsesh@pilani.bits-pilani.ac.in}

\affil[1]{\orgdiv{Department of Computer Science and Information Systems},
\orgname{BITS Pilani, Pilani Campus},
\orgaddress{\street{Vidya Vihar}, \city{Pilani}, \postcode{333031}, \state{Rajasthan}, \country{India}}}

\abstract{Learning-by-teaching tutors built on large language models let students consolidate
their understanding by explaining a topic to an AI learner, but existing systems work only in
English. We show that an English-only tutor can be made multilingual without rebuilding it: a
translation boundary translates only what the participant writes and reads, so inference,
evaluation and stored research data stay in English. We place the same boundary around two
independent Teach2Learn implementations, one with subject-specific fine-tuned learners and one with
a retrieval-augmented GPT-4o learner, and use it to compare them on identical input across all 22
scheduled Indian languages, with 24 reviewed lessons, one independent judge and three
pre-registered protocols. Across 1,575 boundary crossings per system, no text passed untranslated and
no Indian-language script reached the English side. Both systems placed the same 15 languages in the
ready tier and agreed on all 22, and per-language fluency was closely correlated across them,
suggesting that the translation layer is a dominant determinant of multilingual quality. The two
architectures differed in conversational behaviour: the retrieval-augmented learner read slightly
more fluently but repeated a question in 23 of 72 conversations, against none for the fine-tuned
learners, and neither the teacher's earlier sentences nor overlapping retrieval explained the
repeats. \trackb{One sentence on the user study's main result.}}

\keywords{Learning by Teaching, Multilingual Education, Indian Languages, Machine Translation, Large Language Models, Retrieval-Augmented Generation, Fine-Tuning, LLM-as-a-Judge, Speech Input, Mixed-Methods Evaluation}

\maketitle

%% =======================================================================================
\section{Introduction}\label{sec:intro}
%% =======================================================================================

Learning by Teaching (LbT) asks a student to explain a topic to a learner, on the premise that
explaining forces the explainer to organise, justify and repair their own understanding, and that
the learner's questions drive this process \cite{roscoe2007,biswas2005}. Large language models
(LLMs) have made it practical to supply such a learner on demand
\cite{jin2024teach,ma2024debugging,rogers2025playing}. Teach2Learn is one such system, built in two
independent implementations: a multi-agent tool whose student agent is grounded in course textbooks
through retrieval \cite{lahoria2025teach2learn}, and a portal in which subject-specific fine-tuned
models play the learner \cite{finetunedT2L}.

Both implementations, like the other systems cited above, work only in English; the evaluation of
the retrieval-augmented tool names multilingual support as future work \cite{lahoria2025teach2learn}.
Learning in a language one does not fully command is a documented barrier
\cite{unesco2016language}, and India's National Education Policy 2020 makes the home language the
preferred medium of instruction \cite{nep2020}. LbT is an unusually sharp case, because the activity
\emph{is} producing explanations: a student who can explain a CPU scheduler fluently in Tamil but
haltingly in English is being assessed on their English.

A multilingual tutor could replace the learner with a model that works natively in every language,
but that would mean retraining the fine-tuned learners, which are 7--9B-parameter models trained on
English dialogues, and rewriting evaluation logic that parses English output in both
implementations. We take the other route. We draw a single translation boundary around each existing
system: only text the participant writes or reads crosses it, and everything behind it, including
the learner, its evaluation and all stored research data, stays English.

Translating at the edge is not free. Translation can change the technical meaning of an explanation
before the learner reads it, make the learner's questions stilted when they are translated back,
degrade the English that the feedback judge evaluates, and fail without anyone noticing. Whether
these costs depend on the language, on the learner, or on both is an empirical question.

Because both implementations now share the same translation module, we can answer it under
identical conditions. We built a benchmark in which both learners receive the same human-reviewed
teacher scripts, translated by the same code into the same 22 languages and scored by the same
independent judge. The comparison is between two complete learner architectures, not an ablation of
retrieval versus fine-tuning: the learners also differ in base model and prompt. Three protocols were
pre-registered, two of them with predictions.

Four findings stand out. First, the boundary is reliable: across 1,575 crossings per system, no text
passed untranslated and no Indian-language script reached the English side. Second, quality varies
far more by language than by subject or system: both systems place the same 15 languages in the
ready tier, and their per-language fluency is closely correlated (Spearman $\rho = 0.87$). Third,
the two architectures differ in conversational behaviour: the retrieval-augmented GPT-4o learner
reads slightly more fluently, but it repeats a question in one conversation in three, while the
fine-tuned learners never do. Fourth, two pre-registered explanations for those repeats, the
learner not seeing the teacher's earlier sentences and the learner retrieving the same passages
again, were both rejected.

\subsection{Our Contributions}\label{sec:contributions}

\begin{enumerate}
\item A \textbf{translation-boundary architecture} that makes an English-only LbT tutor multilingual
without changing its inference, evaluation or storage, implemented with a byte-identical module in two
independent systems. The architecture also accepts speech input, which joins the same path as typed
text; we describe it but do not evaluate recognition accuracy (Section~\ref{sec:system}).
\item A \textbf{commensurable cross-system benchmark}: a shared, human-reviewed lesson bank, identical
translation code, one judge model and prompt for both systems, mechanical failure checks alongside
judged quality, and three pre-registered protocols (Section~\ref{sec:method}).
\item \textbf{Findings} on multilingual quality (Section~\ref{sec:results}) and on learner behaviour
(Section~\ref{sec:behaviour}), including two pre-registered negative results on why the
retrieval-augmented learner repeats itself.
\item \trackb{A multilingual user study with students and instructors.}
\end{enumerate}

\subsection{Research Questions}\label{sec:rqs}

\begin{description}
\item[RQ1] Can a single translation layer at the edge of the system extend two architecturally
different LbT tutors to 22 Indian languages without changing inference, evaluation or stored data,
and how reliably does it work?
\item[RQ2] How much does multilingual quality vary by language, by subject and by system?
\item[RQ3] Given identical teacher input, how do the two learner architectures differ in fluency,
repetition, responsiveness, lexical variety and question depth?
\item[RQ4] \trackb{How do students experience teaching an AI learner in their own language?}
\end{description}

After related work (Section~\ref{sec:related}) and the systems (Section~\ref{sec:system}), we present
the benchmark (Section~\ref{sec:method}), its results for RQ1--RQ2 (Section~\ref{sec:results}) and RQ3
(Section~\ref{sec:behaviour}), and the user study for RQ4 (Sections~\ref{sec:study-method}
and~\ref{sec:study-findings}), followed by discussion, limitations and conclusions.

%% =======================================================================================
\section{Related Work}\label{sec:related}
%% =======================================================================================

This section reviews the pedagogy behind LbT, LLM-based LbT systems, the two techniques our two
learners use, and the multilingual, speech and evaluation work our design draws on.

\subsection{Learning by Teaching: Pedagogical Foundations}

Tutor learning arises when explaining leads the tutor to build new knowledge rather than merely
restate it. Roscoe and Chi distinguish these as knowledge-building and knowledge-telling, and show
that the tutee's questions are a main driver of the former \cite{roscoe2007}. Teachable-agent
systems operationalise this with a computational tutee. In Betty's Brain, students teach a causal
map to an agent that is then quizzed on it \cite{biswas2005,blair2007}; SimStudent learns
problem-solving skills from students' demonstrations and feedback, and the quality of that tutoring
predicts the students' own gains \cite{matsuda2013}. Both lines locate the benefit in the tutee's
questioning, which is why we measure what the learner asks as well as how it reads after
translation.

\subsection{LLM-based Learning by Teaching Systems}

LLMs let the tutee converse in open language. AlgoBo constrains an LLM's knowledge so that students
can teach it programming \cite{jin2024teach}. Related work uses an LLM tutee for debugging
\cite{ma2024debugging}, deploys a teachable ChatGPT agent in programming courses \cite{chen2024bjet},
studies an LLM that ``plays dumb'' in university CS classes \cite{rogers2025playing}, and embeds LbT
in a conversational tutoring system \cite{schmucker2024ruffle}. Teach2Learn extends this line to
whole CS courses with Bloom-ordered questioning, first with a retrieval-augmented multi-agent design
\cite{lahoria2025teach2learn} and then with fine-tuned learner models \cite{finetunedT2L}. All of these
systems operate in English; in the literature we reviewed, we found none that supports Indian
languages.

\subsection{Retrieval-Augmented and Fine-Tuned Learner Agents}

Retrieval-augmented generation conditions each response on passages retrieved from a document
collection \cite{lewis2020rag}; in a tutor it can confine the learner to course material
\cite{lahoria2025teach2learn}. Fine-tuning changes the model's weights instead, and LoRA makes this
practical for 7--9B-parameter models \cite{hu2022lora}; the fine-tuned Teach2Learn uses it to train
small models to behave as novice learners who always ask \cite{finetunedT2L}. Comparisons of the two
have mostly concerned knowledge injection, where retrieval outperformed unsupervised fine-tuning
\cite{ovadia2024}. Whether either yields a better \emph{learner}, one that asks relevant,
non-repetitive questions, is a different question; in the literature we reviewed we found no study
that compares the two on the same teacher input. Our comparison (Section~\ref{sec:behaviour})
addresses it only in part, because the two learner architectures also differ in base model and
prompt.

\subsection{Multilingual and Indian-language Educational Technology}

LLM quality falls away from English, most steeply for low-resource languages and non-Latin scripts
\cite{ahuja2023mega}. Two designs follow: use a natively multilingual model end to end, or keep an
English pipeline and translate at its edges. For India, IndicTrans2 was the first open translation
model to cover all 22 scheduled languages \cite{gala2023indictrans2}. We use a commercial service,
Sarvam AI's \texttt{sarvam-translate:v1}, which covers the same languages behind one API
\cite{sarvamtranslate}. Section~\ref{sec:boundary} explains why edge translation suits systems whose
learner is an English-trained adapter and whose evaluation logic parses English strings. Real
students also mix English technical terms into Indian-language sentences, a case our scripted
benchmark does not exercise (Section~\ref{sec:lim-benchmark}).

\subsection{Speech Input for Indian Languages}

Whisper made multilingual speech recognition broadly available \cite{radford2023whisper}, but its
training data is thin for many Indian languages. The Vistaar benchmarks document large differences
in word error rate across Indian languages and domains, and show gains from fine-tuning on
Indian-language data \cite{bhogale2023vistaar}. This coverage gap decided our default speech backend
(Section~\ref{sec:voice}).

\subsection{LLMs as Judges of Translation and Dialogue}

LLM judges can rate translation quality at or near the level of established metrics
\cite{kocmi2023gemba} and agree with human preferences on open-ended chat \cite{zheng2023judging}, but
they are sensitive to presentation order and other surface features \cite{wang2024fair}. Calibrated
against 20,000 human judgments in eight languages, GPT-4 evaluators scored too high, most of all in
low-resource and non-Latin-script languages \cite{hada2024multilingual}, which is where our weakest
languages sit. We therefore report mechanical checks that need no judge alongside judged scores,
measure the judge's agreement with itself, and treat judged scores as directional until native
readers confirm them. To characterise the learner's questions we use Graesser and Person's question
taxonomy \cite{graesser1994}.

\textbf{Research gap.} In the literature we reviewed, no LbT tutor supports Indian languages, and
no study compares a fine-tuned learner with a retrieval-augmented one on identical teacher input.

%% =======================================================================================
\section{System Architecture and Learning Workflow}\label{sec:system}
%% =======================================================================================

\subsection{The Two Teach2Learn Systems}\label{sec:two-systems}

Table~\ref{tab:systems} summarises the two systems. In both, the participant chooses a subject and
subtopic and teaches it to an AI learner that asks follow-up questions. The four subjects are BITS
Pilani's first courses in Computer Programming (\texttt{cp}), Data Structures and Algorithms
(\texttt{dsa}), Object-Oriented Programming in Java (\texttt{oop\_java}) and Operating Systems
(\texttt{os}). Below we call the two learner architectures the \emph{fine-tuned learners} and the
\emph{RAG learner}.

\begin{table}[t]
\caption{The two learner architectures compared in this work. Only the translation and speech
modules are shared (Section~\ref{sec:boundary}).}\label{tab:systems}
\footnotesize
\begin{tabular*}{\textwidth}{@{\extracolsep\fill}p{0.2\textwidth}p{0.36\textwidth}p{0.36\textwidth}}
\toprule
& Fine-tuned Teach2Learn \cite{finetunedT2L} & RAG Teach2Learn \cite{lahoria2025teach2learn} \\
\midrule
Learner model & One LoRA adapter per base model \cite{hu2022lora}; subject routes to Qwen3 8B (\texttt{cp}) \cite{qwen3}, Gemma 2 9B (\texttt{dsa}) \cite{gemma2}, Mistral 7B (\texttt{oop\_java}) \cite{mistral7b}, Granite 3.3 8B (\texttt{os}) \cite{granite33}; Qwen 2.5 7B \cite{qwen25} also deployed & GPT-4o \cite{gpt4o} via OpenRouter for every subject \\
Learner context & Full dialogue history in the adapter's training format & Four textbook passages retrieved per teacher turn ($k=4$), plus its own earlier replies \\
Self-evaluation & Private inline block (Bloom level, score, strength, improvement) that drives session completion & Separate evaluator, scorer and summariser agents \\
Stack & FastAPI, SQLite, vLLM \cite{kwon2023vllm} & Chainlit, LangChain, MongoDB, Chroma \\
Corpus & None at inference & 21,564 chunks from 7,093 textbook pages across the four courses (Table~\ref{tab:corpus}) \\
\botrule
\end{tabular*}
\end{table}

\paragraph{Fine-tuned learners.} Each subject routes to its own LoRA-adapted base model, trained on
English dialogues in which the teacher's turn is rendered as \texttt{Previous Question: \ldots\
Teacher's Answer: \ldots}. The learner emits a follow-up question together with a private evaluation
block giving a Bloom level and a Low/Medium/High score. The portal strips that block before display
and uses it to end a session early once a knowledge threshold is reached, between three and six
teacher turns. The training data and procedure are described in \cite{finetunedT2L}.

\paragraph{RAG learner.} The student agent is GPT-4o prompted to behave as a novice while privately
consulting four textbook passages retrieved for each teacher turn (Appendix~\ref{app:prompts}; corpus
in Appendix~\ref{app:implementation}). It returns a rating, a reflection and an optional question;
when it rates an explanation ``understood'' it may return no question, and the participant then sees
the reflection.

\subsection{The Translation Boundary}\label{sec:boundary}

Figure~\ref{fig:boundary} shows the design. Everything that existed before (inference, the private
evaluation, completion logic, the feedback judge and storage) sits inside an English zone that does
not know a language was chosen. Text crosses the boundary twice per turn: the participant's answer is
translated into English on the way in, and the learner's question is translated into the
participant's language on the way out. The opening question, the completion message and the final
feedback are translated outward in the same way (Table~\ref{tab:crosses}).

\begin{figure}[t]
\centering
\includegraphics[width=\textwidth]{figures/fig1-boundary.pdf}
\caption{The translation boundary. Only participant-visible text crosses it; the model, its
evaluation, the judge and all stored data stay English.}\label{fig:boundary}
\end{figure}

\begin{table}[t]
\caption{What crosses the boundary.}\label{tab:crosses}
\footnotesize
\begin{tabular*}{\textwidth}{@{\extracolsep\fill}p{0.36\textwidth}p{0.22\textwidth}p{0.34\textwidth}}
\toprule
Content & Translated? & Reason \\
\midrule
Participant's typed or dictated answer & Into English & The learner and judge are English systems \\
Learner's follow-up question & Into participant's language & Participant-visible \\
Opening question, completion message & Into participant's language & Participant-visible \\
Final feedback (strengths, improvements, summary) & Into participant's language & Participant-visible \\
Private evaluation block & Never & Completion logic parses English labels \\
Judge prompts and judge output & Never & Judge reads the English answer \\
Database rows and transcripts & Never; stored in English & Research analysis stays language-independent \\
\botrule
\end{tabular*}
\end{table}

Three design choices follow from keeping the English zone untouched.

\paragraph{Store English, translate at display time.} The database and transcripts hold only the
English text, and stored conversations are translated again when a page is rendered. A Hindi session
therefore yields an English transcript that existing analysis reads unchanged. The cost is that the
participant's original wording is not kept (Section~\ref{sec:disc-storage}).

\paragraph{The judge sits inside the English zone.} Both systems' feedback evaluators read the
translated-to-English answer, so translation quality feeds evaluation quality. In the fine-tuned
portal, where evaluation decides when a session ends, a translation problem becomes a session-length
problem.

\paragraph{Fail open.} On any translation failure the module returns the original text and the
session continues, so a degraded session looks like a working one; the benchmark therefore checks
every crossing mechanically (Section~\ref{sec:measures}).

\subsection{Implementation and Language Coverage}\label{sec:impl}

Both systems import one translation module, byte-identical in the two code bases, that calls Sarvam
AI's \texttt{sarvam-translate:v1} \cite{sarvamtranslate}. It also rejects one failure that is not an
error: for particular sentence--language pairs the service returned a fragment of its own API schema
in place of a translation, which would otherwise have reached the learner as the participant's words.
The first release offered eight languages; extending it to all 22 languages the service accepts
changed two dictionaries in one file, because the language picker, session metadata and every
display path read from them. The fine-tuned portal translates the conversation and feedback but not
its interface; the RAG tool also translates the interface messages it sends. Appendix~\ref{app:implementation}
gives the details, and Appendix~\ref{app:languages} lists the languages, scripts and codes.

\subsection{Speech Input}\label{sec:voice}

Both systems also accept spoken answers. A second shared module sends the recording to a
speech-to-text service and discards the audio; the returned text then crosses the boundary exactly
like a typed answer, so speech needs no separate path through the English zone. The default service,
Sarvam's Saaras v3, is documented as covering all 22 languages \cite{sarvamstt}; the alternative,
Whisper large-v3 \cite{radford2023whisper}, lists 14 of the 22 \cite{whispertokenizer}. Recognition
accuracy is outside the scope of this study (Appendix~\ref{app:implementation}).

\subsection{Learning Workflow in the Participant's Language}\label{sec:workflow}

A session follows the English-only workflow of the earlier papers \cite{lahoria2025teach2learn,finetunedT2L}
with three changes. The participant first chooses a teaching language, which is stored with the
session. Each explanation, typed or dictated, is translated into English, stored and passed to the
learner (in the RAG tool it is also the retrieval query), and the learner's English question is
translated before it is shown; evaluation and the completion rule run on the English text. Final
feedback is produced in English and translated when displayed.

\subsection{An Example Turn}\label{sec:example}

Figure~\ref{fig:example} traces one Hindi turn from the benchmark (Section~\ref{sec:method}). The
teacher's English sentence is translated into Hindi, standing in for what a participant would type.
The inbound translation that both learners receive keeps the claim but drops ``on the CPU'' and
``runnable''. The fine-tuned learner (Granite 3.3 8B, the operating-systems model) opens with a stock
expression of doubt before asking why the choice matters; the RAG learner asks how the scheduler
chooses among ready processes. Both Hindi questions address the teacher with feminine verb forms, a
choice the translation service made on the participant's behalf (Section~\ref{sec:lim-system}).

\begin{figure}[t]
\centering
\includegraphics[width=\textwidth]{figures/fig2-example-hindi.pdf}
\caption{One Hindi turn as it crosses the boundary (lesson
\texttt{os/threads-scheduling/cpu-scheduling}, turn 1). Blue rows are inside the English zone; grey
rows are what the participant sees. Both systems received the identical Hindi sentence.}\label{fig:example}
\end{figure}

%% =======================================================================================
\section{Benchmark Methodology}\label{sec:method}
%% =======================================================================================

\subsection{Commensurability by Construction}\label{sec:commensurable}

Any number used to compare the two systems must come from a component that is identical for both,
applied to inputs that are identical for both. Three components are shared.

\begin{enumerate}
\item \textbf{Translation.} The byte-identical module of Section~\ref{sec:impl}.
\item \textbf{Teacher input.} One lesson bank, read from the same file by both harnesses
(Section~\ref{sec:lessons}).
\item \textbf{Judge.} One judge model with the same fidelity and fluency prompts
(Section~\ref{sec:measures}).
\end{enumerate}

The comparison is architectural: the learners also differ in base model, size and prompt. Their
native self-assessments are excluded, because the scales do not map onto each other (an inline
three-level grade against five continuous sub-scores from separate agents). Two inputs cannot be
made identical without disabling part of a system. The opening question comes from the adapter on one
side and from a retrieval-grounded generator on the other, so question measures use follow-ups only.
The RAG student may reply with a reflection instead of a question, whereas the adapter always asks,
so the share of turns that are questions is reported but not scored (Table~\ref{tab:notscored}).

\subsection{Lesson Bank}\label{sec:lessons}

A teacher is simulated by a fixed script, so that differences between runs reflect the language or
the system rather than what a participant happened to say. The bank holds 24 lessons (four subjects,
three topics each, two subtopics per topic) drawn from the portal's own catalogue. Each lesson is three
sentences, each a claim plus its reason, with operators spelled out in words because every sentence is
machine-translated (Appendix~\ref{app:lessons}). A reader who knows the material reviewed every script
before any run, and the harness refuses to run an unreviewed one.

\subsection{Runs}\label{sec:runs}

\begin{description}
\item[Sweep A (English).] All 24 lessons, three repeats each, per system (72 conversations), without
translation or judge.
\item[Sweep B (languages).] Eight lessons, two per subject, in each of 22 languages, one repeat per
system: 176 language-runs of three turns.
\item[Model run.] All five deployed fine-tuned models on the \texttt{os} and \texttt{cp} lessons in
English (Appendix~\ref{app:models}).
\item[Pre-registered analyses.] The teacher-history experiment (Sweep A with one change) and the
retrieval-overlap analysis (retrieval only, no generation), both on the RAG side.
\end{description}

In a Sweep B turn the English teacher sentence is first translated into the target language (the
\emph{seed}, standing in for what a participant would type), then translated back to English
(\emph{inbound}), passed to the learner, and the learner's English reply is translated out
(\emph{outbound}). Each turn therefore makes three crossings, and 525 turns make 1,575. The opening
question's translation is checked too but not counted in that total.

Both systems sampled with top-$p$ 0.9. The RAG learner (GPT-4o, pinned) used temperature 0.25 as
specified and the fine-tuned runs 0.2 (Section~\ref{sec:deviations}). Every results file records its
full configuration.

\subsection{Measures}\label{sec:measures}

\paragraph{Mechanical checks.} These need no judgement: whether a seed or outbound translation came
back unchanged (untranslated), whether any Indian-language character reached the English side (a
script leak), how many retries a crossing needed, and whether a schema payload was rejected.

\paragraph{Judged quality.} An independent model, \texttt{openai/gpt-5.6-luna} accessed through
OpenRouter \cite{gptluna}, scores two things on a 1--5 scale with the prompts in
Appendix~\ref{app:prompts}. \emph{Fidelity} compares the English teacher sentence with its seed
translation: did the technical meaning survive, in the right script? \emph{Fluency} reads the
learner's outbound message alone: is it natural text a learner would write? Fidelity grades the
translation service on input that is identical for both systems, so it serves as a control rather
than a comparison. Fluency is the judged measure that depends on the learner.

\paragraph{Language tiers.} A language is \emph{Blocked} if its mean fidelity is below 2.0,
\emph{Review} if mean fidelity is below 3.5 or mean fluency below 4.0, and \emph{Ready} otherwise.

\paragraph{Conversation measures (Sweep A).} For each follow-up we take its question sentence (the
last sentence ending in a question mark, which drops stock lead-ins such as ``That makes sense,
but\ldots''). A conversation \emph{repeats} if the same normalised question sentence occurs twice.
\emph{Lexical variety} is the entropy of question openers (first six words), normalised by
$\log n$ and averaged over 400 random subsamples of equal size for both systems. As a crude proxy for
responsiveness, a question \emph{answers the current sentence} if it shares the most content words
uniquely with the teacher's latest sentence, and \emph{looks backward} if an earlier sentence shares
more.

\paragraph{Question depth.} Each question sentence is sent alone, with no system name, subject or
teacher text and with all arms shuffled together, to the same judge, which assigns one of Graesser and
Person's 18 categories \cite{graesser1994}. Code, not the judge, groups the categories into shallow,
intermediate and deep levels, as pre-registered (Appendix~\ref{app:prereg}). A random 60 questions are
judged twice to measure agreement.

\subsection{Behavioural and Translation Dimensions}\label{sec:scorecard}

RQ3 is answered along six dimensions, each placed on a 0--100 scale and reported separately
(Table~\ref{tab:dimensions}): fluency, $(\bar{f}-1)/4\times100$ over paired turns; question depth, the
share of questions rated deep; lexical variety, the normalised opener entropy; non-repetition, the
share of English conversations without a repeated question; robustness, the share of the 22 languages
rated Ready; and reliability, the share of crossings with no retry, untranslated result or leak. Only
signals that both harnesses record and that the learner itself produces are used. A dimension favours
a system only when the 95\% interval of the difference excludes zero; for reliability we use Fisher's
exact test, and robustness is a tie unless the Ready counts differ. We do not combine the dimensions
into an overall score in the main text, because any weighting is necessarily normative; the
pre-registered composite and its sensitivity to weighting are reported in Appendix~\ref{app:composite}.

\subsection{Statistical Analysis}\label{sec:stats}

Fluency differences use turns paired by (lesson, language, turn) and a 5,000-draw cluster bootstrap
over (lesson, language) \cite{efron1993}. Resampling lessons with replacement duplicates questions,
which depresses entropy below its own point estimate, so the English measures instead draw half the
lessons without replacement (1,000 draws) and centre each distribution on the full-sample estimate.
Draws are paired across systems throughout. Agreement between the two systems' per-language means is
summarised by Pearson and Spearman correlations. The retrieval-overlap analysis uses a 5,000-draw
bootstrap that resamples lessons. All analyses use fixed random seeds. The analysis scripts and raw
results files are listed under Code and Data availability.

\subsection{Pre-Registration}\label{sec:prereg}

Three protocols were written down, with UTC timestamps, before the data they judge existed
\cite{nosek2018}; the full text is in Appendix~\ref{app:prereg}.

\begin{enumerate}
\item \textbf{Teacher-history experiment} (2026-09-14 18:31 UTC): give the RAG student the teacher's
earlier sentences and change nothing else; predicted to reduce conversations that repeat a question
from 23 of 72 toward the fine-tuned level of 0.
\item \textbf{Question-depth instrument} (2026-09-14 18:35 UTC): the category grouping, and a rule that
depth is not used if re-judge agreement on depth level falls below 70\%.
\item \textbf{Retrieval-overlap analysis} (2026-09-19 05:32 UTC): predicted that pairs of turns that
retrieve overlapping passages repeat a question more often than pairs that do not.
\end{enumerate}

For both predictions, the failure criterion is that the 95\% interval of the predicted effect
includes zero.

\subsection{Deviations and Revisions}\label{sec:deviations}

Four departures from plan are recorded rather than corrected. (i)~The fine-tuned runs used
temperature 0.2 instead of 0.25; a lower temperature should, if anything, favour repetition on that
side, which makes its non-repetition result conservative. (ii)~The two harnesses first used different
fluency prompts; the fine-tuned side's 525 fluency scores were re-judged with the shared prompt
(Appendix~\ref{app:prompts}) without regenerating anything, and 120 scores moved, almost all by one
point, with no language changing tier. (iii)~The composite changed twice; both earlier versions are in
Appendix~\ref{app:composite}. (iv)~Rate limiting during the fine-tuned side's Sweep B returned text
unchanged for three languages; the harness stopped rather than record clean-looking rows, and 19 cells
were re-run.

\subsection{Cost}\label{sec:cost}

Costs were measured from the provider's credit balance before and after each run, not estimated.
The RAG side's benchmark runs cost US\$3.44 in OpenRouter calls (generation and judge) and INR~430.67
(about US\$5.15) in translation, about US\$8.59 in total. On the fine-tuned side, where generation runs
on local GPUs, 1,060 judge calls cost about US\$0.30. The teacher-history experiment cost US\$0.89, and
the retrieval-overlap analysis made 144 embedding calls (two identical runs) at a cost of a fraction
of a cent.

\subsection{Use of AI Tools}\label{sec:ai-use}

This subsection reports AI assistance in preparing the software and this article, as Springer
Nature's policy requires. It is separate from the language models studied in this paper, which are
described above: the learners, the judge, and the translation and speech services.

\paragraph{Software.} Claude Opus 5 and Claude Sonnet 5 (Anthropic) assisted in writing the
translation and transcription modules, the benchmark harnesses, and the analysis and reporting
scripts. Both were used through the Claude Code command-line tool between 18 August and 19 September
2026. The authors directed each change, reviewed it, and ran the test suites. Commits made with AI
assistance are marked as co-authored in the project's version history.

\paragraph{Manuscript.} Claude Opus 5 (model identifier \texttt{claude-opus-5}), used through Claude
Code between 16 and 19 September 2026, wrote the first draft of Sections~\ref{sec:intro}--\ref{sec:behaviour},
\ref{sec:discussion}, \ref{sec:lim-system}, \ref{sec:lim-benchmark} and~\ref{sec:conclusion} and of
the appendices, and revised them after an internal review. It worked from the project's own results
files, documentation and code, following an outline agreed with the authors. It also generated the
long tables and the figures by writing scripts that read the committed results files; the figures are
data plots, and no generative image model was used. It proposed references, whose bibliographic
details were checked against publisher or indexing records. The authors' instructions to the tool are
available from the corresponding author on request. The authors reviewed, revised and verified the
text, numbers and references, and take full responsibility for the content. No AI tool is listed as
an author.
\todo{Before submission, confirm that every statement in this subsection is true, and add how the
user-study sections were written.}

%% =======================================================================================
\section{Results: The Multilingual Boundary}\label{sec:results}
%% =======================================================================================

\subsection{Reliability of the Boundary (RQ1)}\label{sec:res-boundary}

Table~\ref{tab:mechanical} gives the mechanical checks. Across 525 judged turns per system, no
participant text and no text shown back was left untranslated, and no Indian-language script reached
the English zone. Four of the fine-tuned side's 1,575 crossings needed a retry, all rate-limit
responses from the translation service during that run; the RAG side needed none. The difference is
not significant (Fisher's exact $p=0.12$) and does not involve the learner.

\begin{table}[t]
\caption{Mechanical checks over Sweep B. No judgement involved.}\label{tab:mechanical}
\footnotesize
\begin{tabular*}{\textwidth}{@{\extracolsep\fill}lrrr}
\toprule
Check & Fine-tuned & RAG & Over (per system) \\
\midrule
Language-runs completed & 175 & 175 & 176 \\
Untranslated participant text & 0 & 0 & 525 turns \\
Untranslated text shown back & 0 & 0 & 525 turns \\
Native script in the English zone & 0 & 0 & 525 turns \\
Crossings needing a retry & 4 & 0 & 1,575 crossings \\
Schema payloads rejected & 1 & 1 & 1,575 crossings \\
\botrule
\end{tabular*}
\end{table}

Two defects surfaced. One sentence of the dynamic-memory lesson, translated into Punjabi, returned the
service's API schema on every attempt on both systems; the guard of Section~\ref{sec:impl} rejected
it, so it reached neither learner, and Punjabi's other lessons ran normally. The second defect is
upstream and unfixed: Sindhi is returned in Devanagari, whereas Perso-Arabic is the dominant script
for written Sindhi, and the service does not accept an output-script parameter for this model. Sindhi
scored fidelity between 1.0 and 1.5 on all four fine-tuned models and on the RAG learner.

\subsection{Language Dominates Subject and System (RQ2)}\label{sec:res-languages}

Table~\ref{tab:languages} gives fidelity and fluency per language for both systems, and
Appendix~\ref{app:matrices} breaks each down by subject. Both systems place 15 languages in the Ready
tier, 6 in Review and 1 in Blocked, and \textbf{they agree on the tier of every one of the 22
languages}. The Review languages (Bodo, Kashmiri, Konkani, Manipuri, Santali and Sanskrit) keep high
fidelity but fall short on fluency; Sanskrit is the weakest of them on both systems, and Sindhi is
Blocked on both.

\begin{table}[p]
\caption{Fidelity and fluency by language, 1--5, from one judge and one prompt for both systems.
$\Delta$ is RAG minus fine-tuned (FT). Tiers agree for all 22 languages.}\label{tab:languages}
\footnotesize
\begin{tabular*}{\textwidth}{@{\extracolsep\fill}llrrrrrrl}
\toprule
& & \multicolumn{3}{c}{Fidelity} & \multicolumn{3}{c}{Fluency} & \\
\cmidrule(lr){3-5}\cmidrule(lr){6-8}
Language & Code & FT & RAG & $\Delta$ & FT & RAG & $\Delta$ & Tier \\
\midrule
Punjabi & \texttt{pa} & 4.67 & 4.57 & $-$0.10 & 4.76 & 4.95 & $+$0.19 & Ready \\
Malayalam & \texttt{ml} & 4.75 & 4.83 & $+$0.08 & 4.58 & 4.92 & $+$0.33 & Ready \\
Telugu & \texttt{te} & 4.50 & 4.33 & $-$0.17 & 4.79 & 4.83 & $+$0.04 & Ready \\
Urdu & \texttt{ur} & 4.67 & 4.67 & $+$0.00 & 4.54 & 4.83 & $+$0.29 & Ready \\
Gujarati & \texttt{gu} & 4.71 & 4.71 & $+$0.00 & 4.75 & 4.79 & $+$0.04 & Ready \\
Hindi & \texttt{hi} & 4.58 & 4.62 & $+$0.04 & 4.62 & 4.79 & $+$0.17 & Ready \\
Bengali & \texttt{bn} & 4.38 & 4.38 & $+$0.00 & 4.79 & 4.71 & $-$0.08 & Ready \\
Odia & \texttt{or} & 4.75 & 4.54 & $-$0.21 & 4.33 & 4.71 & $+$0.38 & Ready \\
Nepali & \texttt{ne} & 4.42 & 4.42 & $+$0.00 & 4.58 & 4.71 & $+$0.12 & Ready \\
Kannada & \texttt{kn} & 4.50 & 4.54 & $+$0.04 & 4.50 & 4.67 & $+$0.17 & Ready \\
Marathi & \texttt{mr} & 4.46 & 4.38 & $-$0.08 & 4.54 & 4.67 & $+$0.12 & Ready \\
Assamese & \texttt{as} & 4.21 & 4.54 & $+$0.33 & 4.58 & 4.62 & $+$0.04 & Ready \\
Tamil & \texttt{ta} & 4.50 & 4.38 & $-$0.12 & 4.58 & 4.58 & $+$0.00 & Ready \\
Maithili & \texttt{mai} & 4.62 & 4.58 & $-$0.04 & 4.33 & 4.38 & $+$0.04 & Ready \\
Dogri & \texttt{doi} & 4.38 & 4.42 & $+$0.04 & 4.00 & 4.25 & $+$0.25 & Ready \\
Kashmiri & \texttt{ks} & 4.12 & 4.50 & $+$0.38 & 3.79 & 3.96 & $+$0.17 & Review \\
Bodo & \texttt{brx} & 4.75 & 4.75 & $+$0.00 & 3.79 & 3.96 & $+$0.17 & Review \\
Konkani & \texttt{kok} & 4.58 & 4.58 & $+$0.00 & 3.38 & 3.83 & $+$0.46 & Review \\
Manipuri & \texttt{mni} & 4.79 & 4.88 & $+$0.08 & 3.67 & 3.79 & $+$0.12 & Review \\
Santali & \texttt{sat} & 4.75 & 4.88 & $+$0.12 & 3.42 & 3.57 & $+$0.15 & Review \\
Sanskrit & \texttt{sa} & 3.17 & 3.25 & $+$0.08 & 3.00 & 3.25 & $+$0.25 & Review \\
Sindhi & \texttt{sd} & 1.33 & 1.00 & $-$0.33 & 2.12 & 1.88 & $-$0.25 & Blocked \\
\botrule
\end{tabular*}
\end{table}

Agreement on fidelity is expected by construction, since fidelity grades the translation of the same
teacher sentence on both sides. The informative measure is fluency, which grades each learner's own
message after translation. Here too the two systems agree closely (Figure~\ref{fig:agreement}): across
the 22 languages, the correlation between their mean fluency is $r = 0.98$ (Pearson) and $\rho = 0.87$
(Spearman), and $\rho = 0.86$ without Sindhi. Every language falls on the same side of the 4.0 Review
threshold in both systems (Dogri sits exactly on it for the fine-tuned learners), and the six Review
languages plus Sindhi are the seven lowest on both.

\begin{figure}[t]
\centering
\includegraphics[width=0.62\textwidth]{figures/fig7-language-agreement.pdf}
\caption{Mean judged fluency per language, fine-tuned learners against the RAG learner. Points near
the diagonal are languages the two architectures handle equally well.}\label{fig:agreement}
\end{figure}

Language accounts for far more of the variation than subject (Table~\ref{tab:subjects}). Mean fluency
spans 0.27 across subjects on the fine-tuned side and 0.12 on the RAG side, but 2.67 and 3.07 across
languages. The two learners share nothing but the translation layer, so this pattern is consistent
with the translation layer being a dominant determinant of multilingual quality in this
architecture, and with technical vocabulary not being the binding constraint. It does not isolate
the translation layer's contribution: observed fluency combines inbound translation, the learner's
English output, outbound translation and the judge.

\begin{table}[t]
\caption{Quality by subject, 1--5. Subject also selects the fine-tuned model.}\label{tab:subjects}
\footnotesize
\begin{tabular*}{\textwidth}{@{\extracolsep\fill}llrrrr}
\toprule
& & \multicolumn{2}{c}{Fidelity} & \multicolumn{2}{c}{Fluency} \\
\cmidrule(lr){3-4}\cmidrule(lr){5-6}
Subject & FT model & FT & RAG & FT & RAG \\
\midrule
Computer Programming & Qwen3 8B & 4.38 & 4.35 & 4.31 & 4.31 \\
Data Structures and Algorithms & Gemma 2 9B & 4.35 & 4.36 & 4.14 & 4.31 \\
OOP with Java & Mistral 7B & 4.34 & 4.39 & 4.04 & 4.23 \\
Operating Systems & Granite 3.3 8B & 4.30 & 4.30 & 4.13 & 4.35 \\
\midrule
Range across subjects & & 0.08 & 0.09 & 0.27 & 0.12 \\
\botrule
\end{tabular*}
\end{table}

Two properties of the measurement qualify these numbers. The translation service is not
deterministic: 116 of the 525 teacher sentences were translated differently in the two systems'
runs, concentrated in low-resource languages (Santali 14 of 24, Kashmiri 12 of 24). And the judge is
noisy at the scale of the differences we report: on the 409 turns where both systems received
identical translated sentences, it gave identical fidelity scores 84\% of the time, its mean absolute
disagreement with itself on the same sentence was 0.18, and the net fidelity gap was $-0.02$, which we
take as the noise floor.

\subsection{Fluency of the Two Learners after Translation}\label{sec:res-fluency}

The RAG learner's messages read slightly more fluently after translation: $+0.14$ on the 1--5 scale
over 524 paired turns, 95\% CI [$+0.08$, $+0.21$]. The effect is smaller than the judge's disagreement
with itself, but it is consistent: the RAG message scored higher on 147 turns and the fine-tuned one
on 82 (295 level), and RAG was higher in 19 of 22 languages. It survives every cut in
Table~\ref{tab:fluency}, including turns on which the translation service returned identical input to
both systems and pairs matched for length. Computer Programming, where the fine-tuned side runs
Qwen3, is the exception with no difference (Figure~\ref{fig:forest}).

\begin{table}[t]
\caption{Fluency, RAG minus fine-tuned, with 95\% cluster-bootstrap intervals over (lesson,
language).}\label{tab:fluency}
\footnotesize
\begin{tabular*}{\textwidth}{@{\extracolsep\fill}lrrrrl}
\toprule
Cut & Paired turns & FT & RAG & $\Delta$ & 95\% CI \\
\midrule
All paired turns & 524 & 4.16 & 4.30 & $+0.14$ & [$+0.08$, $+0.21$] \\
Questions only & 503 & 4.15 & 4.32 & $+0.17$ & [$+0.11$, $+0.24$] \\
Identical translated input & 409 & 4.20 & 4.34 & $+0.15$ & [$+0.08$, $+0.22$] \\
Length-matched pairs & 164 & 4.05 & 4.20 & $+0.15$ & [$+0.02$, $+0.27$] \\
Computer Programming & 129 & 4.31 & 4.31 & 0.00 & [$-0.13$, $+0.13$] \\
Data Structures and Algorithms & 131 & 4.15 & 4.31 & $+0.16$ & [$+0.02$, $+0.30$] \\
OOP with Java & 132 & 4.04 & 4.23 & $+0.19$ & [$+0.06$, $+0.33$] \\
Operating Systems & 132 & 4.13 & 4.35 & $+0.22$ & [$+0.11$, $+0.33$] \\
\botrule
\end{tabular*}
\end{table}

\begin{figure}[t]
\centering
\includegraphics[width=0.9\textwidth]{figures/fig4-fluency-forest.pdf}
\caption{Fluency difference, RAG minus fine-tuned, with 95\% cluster-bootstrap intervals, for the
cuts in Table~\ref{tab:fluency}.}\label{fig:forest}
\end{figure}

%% =======================================================================================
\section{Results: Learner Behaviour}\label{sec:behaviour}
%% =======================================================================================

This section compares the two learner architectures on the English sweep and, where noted, the
language sweep. Table~\ref{tab:dimensions} summarises the six dimensions. The clearest difference, and
the one that needs no judge, is repetition.

\begin{table}[t]
\caption{The six dimensions, each on 0--100 and reported separately. Intervals are 95\% bootstrap
intervals of the difference; a dimension favours a system only if its interval excludes zero. The
composite that combines them is in Appendix~\ref{app:composite}. $^\dagger$Read with
Section~\ref{sec:res-depth}.}\label{tab:dimensions}
\footnotesize
\begin{tabular*}{\textwidth}{@{\extracolsep\fill}lrrll}
\toprule
Dimension & FT & RAG & Difference [95\% CI] & Favours \\
\midrule
Non-repetition & 100.0 & 68.1 & 0 vs 23 of 72 repeat [$+19$, $+44$ pts] & FT \\
Fluency & 78.9 & 82.5 & $+0.14$ on 1--5 [$+0.08$, $+0.21$] & RAG \\
Lexical variety & 52.9 & 78.5 & $+0.26$ entropy [$+0.14$, $+0.37$] & RAG \\
Question depth$^\dagger$ & 60.2 & 89.9 & $+30$ pts deep [$+18$, $+41$] & RAG \\
Robustness & 68.2 & 68.2 & 15 vs 15 of 22 Ready & Tie \\
Reliability & 99.7 & 100.0 & 4 vs 0 of 1,575, $p=0.12$ & Tie \\
\botrule
\end{tabular*}
\end{table}

\subsection{Repetition and Responsiveness}\label{sec:res-repeats}

The fine-tuned learners never asked the same question twice in a conversation. The RAG learner did so
in 23 of 72 English conversations and in 55 of 175 language conversations, against 5 of 175 for the
fine-tuned learners (Table~\ref{tab:repeats}, Figure~\ref{fig:repeats}). A word-overlap proxy points
the same way: 61\% of the fine-tuned learners' questions relate most closely to the teacher's latest
sentence, against 44\% for RAG, and 30\% of RAG's questions look back to an earlier sentence, against
14\%. The proxy is crude; the repeat count is not. Table~\ref{tab:transcripts} shows one lesson
verbatim. For LbT this matters directly: a learner that re-asks a question the teacher has moved past
gives the teacher nothing new to explain.

\begin{table}[t]
\caption{Repetition and responsiveness. Intervals are for RAG minus fine-tuned.}\label{tab:repeats}
\footnotesize
\begin{tabular*}{\textwidth}{@{\extracolsep\fill}lrrl}
\toprule
Measure & FT & RAG & 95\% CI (pts) \\
\midrule
Conversations repeating a question, English & 0/72 & 23/72 & [$+19$, $+44$] \\
Conversations repeating a question, 22 languages & 5/175 & 55/175 & [$+22$, $+35$] \\
Questions answering the current sentence & 61\% & 44\% & [$-27$, $-7$] \\
Questions looking back to an earlier sentence & 14\% & 30\% & [$+6$, $+28$] \\
\midrule
\multicolumn{4}{l}{Repeating English conversations by subject (of 18): FT 0, 0, 0, 0; RAG 8, 4, 6, 5 (\texttt{cp}, \texttt{dsa}, \texttt{oop\_java}, \texttt{os})} \\
\botrule
\end{tabular*}
\end{table}

\begin{figure}[t]
\centering
\includegraphics[width=0.8\textwidth]{figures/fig5-repeats.pdf}
\caption{English conversations in which the learner asked the same question twice, of 18 per
subject, for the fine-tuned learner, the RAG learner, and the RAG learner given the teacher's
earlier sentences (the teacher-history experiment, Section~\ref{sec:res-explain}).}\label{fig:repeats}
\end{figure}

\begin{table}[t]
\caption{One lesson as each learner asked it (\texttt{os/threads-scheduling/cpu-scheduling}).
Stock lead-ins of the fine-tuned learner are in italics. Bold marks a question asked word for word
earlier in the same conversation.}\label{tab:transcripts}
\footnotesize
\begin{tabular*}{\textwidth}{@{\extracolsep\fill}p{0.04\textwidth}p{0.28\textwidth}p{0.3\textwidth}p{0.3\textwidth}}
\toprule
Turn & Teacher & Fine-tuned (Granite 3.3 8B) & RAG (GPT-4o) \\
\midrule
1 & The scheduler decides which ready process runs next on the CPU, because there are usually more runnable processes than cores. & \emph{That makes sense, but I'm still slightly unsure about one thing.} Why is it important for the scheduler to decide which process runs next? & Could you explain how the scheduler decides which process to run if multiple processes are ready at the same time? \\
2 & Round-robin gives each process a fixed time quantum in turn, so no single process can starve the others while it holds the CPU. & \emph{I think I understand most of that.} How does round-robin scheduling affect system efficiency compared to other methods? & \textbf{Could you explain how the scheduler decides which process to run if multiple processes are ready at the same time?} \\
3 & A very small quantum wastes time on context switches, while a very large one makes the scheduler behave like plain first-come first-served. & \emph{Okay, I'm following you.} Can you give an example of when round-robin might not be the best choice? & Could you explain how the scheduler chooses between multiple ready processes? \\
\botrule
\end{tabular*}
\end{table}

\subsection{Two Pre-Registered Explanations, Both Rejected}\label{sec:res-explain}

\paragraph{The learner cannot see what the teacher said earlier.} The live RAG tool passes the
student only its own earlier replies, not the teacher's earlier sentences. Adding those sentences to
the prompt did not stop the repeats (Appendix~\ref{app:experiment}): 24 of 72 conversations repeated a
question, against 23 of 72 before, a change of $+1.4$ points, 95\% CI [$-10.8$, $+14.2$]. By the
pre-registered criterion the prediction failed, and no secondary measure moved. Subject by subject the
counts shifted in both directions (\texttt{cp} 8$\rightarrow$12, \texttt{oop\_java} 6$\rightarrow$4,
\texttt{os} 5$\rightarrow$4, \texttt{dsa} unchanged at 4), consistent with noise.
\paragraph{The learner retrieves the same passages again.} The student prompt tells the RAG learner to
ask about any gap between the teacher and the retrieved passages, so two turns that retrieve the same
passages might produce the same question. They do not (Table~\ref{tab:retrieval}). No two turns of any
lesson retrieved the same four passages; 57 of the 72 turn pairs shared none. Pairs that shared a
passage repeated a question \emph{less} often than pairs that shared none (5.6\% against 15.8\%, a
difference of $-10.2$ points, 95\% CI [$-18.3$, $-0.8$]), and 54 of the 59 repeated questions came
between turns that retrieved entirely different passages. The pre-registered prediction failed. The
interval lies just below zero and only 15 pairs shared any passage, so we read this as no support for
the hypothesis, not as evidence that shared passages prevent repeats. (Retrieval was deterministic,
and the analysis reproduces the repeat counts above.)

\begin{table}[t]
\caption{Retrieval-overlap analysis (pre-registered). A unit is a lesson and a pair of its turns;
the repeat rate is the share of that lesson's conversations in which the later question repeats the
earlier one. Intervals resample lessons.}\label{tab:retrieval}
\footnotesize
\begin{tabular*}{\textwidth}{@{\extracolsep\fill}lrrrl}
\toprule
& \multicolumn{2}{c}{Repeat rate} & & \\
\cmidrule(lr){2-3}
Arm & Shared passage & None shared & Difference & 95\% CI \\
\midrule
Both arms (primary) & 5.6\% & 15.8\% & $-10.2$ pts & [$-18.3$, $-0.8$] \\
RAG (committed run) & 2.2\% & 16.4\% & $-14.2$ pts & [$-22.6$, $-5.1$] \\
RAG + teacher history & 8.9\% & 15.2\% & $-6.3$ pts & [$-19.6$, $+9.0$] \\
\midrule
\multicolumn{5}{l}{Units: 15 share at least one of four passages, 57 share none; no unit retrieves identical sets.} \\
\multicolumn{5}{l}{Rank correlation of overlap with repeat rate (both arms): $-0.25$ [$-0.47$, $-0.00$].} \\
\botrule
\end{tabular*}
\end{table}

With both explanations rejected, the repeats most likely arise on the generation side. Candidates
include the question template that every RAG question follows (``Could you explain how \ldots?'', see
Section~\ref{sec:res-depth}), the prompt's instruction to assume the explanation was sufficient, and
how the learner uses its memory of its own earlier replies. Testing them requires new generations
with a changed prompt.

\subsection{Lexical Variety}\label{sec:res-lexical}

The fine-tuned learners lean on stock lead-ins. The most common six-word opening begins 83\% of
Qwen3's questions on \texttt{cp} (``that makes sense but i'm still\ldots''), 37\% of Gemma~2's,
26\% of Mistral's and 9\% of Granite's; for the RAG learner the corresponding shares are 15--24\%.
Rarefied to 169 questions per system, normalised opener entropy is 0.53 for the fine-tuned learners
and 0.78 for RAG ($+0.26$ [$+0.14$, $+0.37$]). The measure understates RAG's own formula: its
characteristic opening, ``Could you explain how'', is only four words long, so the six-word opener
still varies. A run of all five fine-tuned models on the same lessons (Appendix~\ref{app:models})
shows that the stock lead-ins depend on the subject as well as the model.

\subsection{Question Depth}\label{sec:res-depth}

The blind judge rated 90\% of the RAG learner's questions deep, against 60\% for the fine-tuned
learners (Table~\ref{tab:depth}, Figure~\ref{fig:depth}). The instrument is consistent: on 60
re-judged questions it assigned the same depth level 97\% of the time and the same category 93\%,
well above the pre-registered 70\% floor.

\begin{table}[t]
\caption{Question depth by arm, from a judge blind to system, subject and lesson.}\label{tab:depth}
\footnotesize
\begin{tabular*}{\textwidth}{@{\extracolsep\fill}p{0.42\textwidth}rrr}
\toprule
& FT & RAG & RAG + teacher history \\
\midrule
Questions judged & 216 & 169 & 177 \\
Deep / intermediate / shallow & 60 / 34 / 6\% & 90 / 7 / 3\% & 92 / 5 / 3\% \\
Deep, each distinct question once per conversation & 60\% of 216 & 90\% of 143 & 91\% of 150 \\
Phrased as a request to explain (``Could you explain/clarify\ldots'') & 20\% & 100\% & 100\% \\
Deep among those / among other phrasings & 66\% / 59\% & 90\% / --- & 92\% / --- \\
Share in the most common category (instrumental/procedural) & 29\% & 82\% & 79\% \\
Categories used, of 18 & 13 & 7 & 7 \\
\botrule
\end{tabular*}
\end{table}

\begin{figure}[t]
\centering
\includegraphics[width=0.85\textwidth]{figures/fig6-depth.pdf}
\caption{Question depth by arm, from a judge blind to system, subject and lesson.}\label{fig:depth}
\end{figure}

We do not think this measures depth alone. \textbf{Every} RAG question is phrased as a request to
explain or clarify, and Graesser and Person file ``how does X work?'' under instrumental/procedural,
a deep category; 82\% of RAG's questions land in that one category, and RAG uses 7 of the 18
categories against 13 for the fine-tuned learners. The judge sees only the question sentence, so it
cannot tell whether the teacher had already answered it. Two checks cut the other way: among the
fine-tuned learners' own questions, those phrased as a request to explain were only slightly more
often deep (66\%) than the rest (59\%), and counting each repeated question once leaves RAG at 90\%.
Neither check can settle RAG's case, because RAG never phrases a question any other way. We therefore
do not claim that the RAG learner asks deeper questions. A person rating a sample of questions with
the lesson in view could separate RAG's template from what its questions ask.

\subsection{Recorded, Not Scored}\label{sec:res-notscored}

Some differences follow from how each system was built or prompted, and scoring them would compare
design choices rather than learners: for example, 100\% of the fine-tuned learners' follow-ups are
questions against 78\% for RAG, whose prompt permits a reflection (Appendix~\ref{app:notscored}).
%% =======================================================================================
\section{User Study: Methodology}\label{sec:study-method}
%% =======================================================================================

\trackb{To be written by the study team. Headings follow the original structure.}

\subsection{Participant Recruitment}

\subsection{Survey - Quantitative Data Collection}

\subsection{Interview - Qualitative Data Collection}

%% =======================================================================================
\section{User Study: Findings}\label{sec:study-findings}
%% =======================================================================================

\trackb{To be written by the study team.}

\subsection{Quantitative Analysis: Students}

\subsection{Quantitative Analysis: Instructors}

\subsection{Qualitative Analysis: Thematic Insights from Students}

\subsection{Qualitative Analysis: Thematic Insights from Instructors}

%% =======================================================================================
\section{Discussion}\label{sec:discussion}
%% =======================================================================================

\subsection{Language Is a Dominant Determinant of Multilingual Quality}

The two systems share only the translation module, yet they place all 22 languages in the same tier,
their per-language fluency is closely correlated, and language explains an order of magnitude more of
the variation than subject. This does not prove that the translation layer alone sets quality, since
inbound translation, the learner and the judge all contribute to every score, but it suggests a
practical priority: when extending a tutor this way, choosing and monitoring the translation service
matters more than choosing the learner. It also offers a cheap screen for a new language: route an
existing English system through the boundary and read the tiers before paying for a natively
multilingual learner.

The same result bounds what edge translation delivers today. The six Review languages are where a
participant is most likely to read stilted questions and where the judge is least trustworthy
\cite{hada2024multilingual}, and Sindhi fails for a reason no application can fix. A deployment should
offer the Ready languages by default, mark the Review languages as experimental, and not offer
Sindhi.

\subsection{An Architectural Comparison, Not an Ablation}

Once language is held constant, the two learners differ in ways that follow their construction. The
RAG learner, a far larger model, writes slightly more fluent messages and varies its openings more.
The fine-tuned learners, trained to always ask and to follow the teacher's latest sentence, never
repeat a question. Neither difference can be credited to retrieval or to fine-tuning as such: the RAG
learner is GPT-4o, the fine-tuned learners are 7--9B-parameter models, and the prompts differ. The
two architectures do not rank on a single scale; they fail in different ways. The apparent depth
advantage of the RAG learner should be read the same way, since a judge that sees only the question
sentence rewards ``Could you explain how \ldots?'', and every RAG question takes that form.

\subsection{Repetition Is the Failure That Matters}

A repeated question is easy to dismiss as a stylistic flaw, but LbT depends on it not happening.
Tutor learning comes from questions that push the tutor to build on what they have just said
\cite{roscoe2007}. A learner that returns to a question the teacher has already moved past gives the
teacher nothing new to explain, and one conversation in three did so with the RAG learner. Two
natural explanations, missing teacher history and repeated retrieval, were each pre-registered and
each rejected. That narrows the search to how the learner generates its questions, which is where
the next experiment should look.

\subsection{A Boundary That Fails Open Must Fail Loudly}

Passing the original text through on failure keeps sessions alive during outages, but a failure
then looks like success. Two defects we found would have been invisible without mechanical checks: a
schema fragment returned in place of a Punjabi translation, and a rate limit that returned text
unchanged. The general principle is that a boundary that fails open must fail loudly. Deployments
should log every passthrough, run the benchmark's checks (retries, schema rejection, script detection)
on the live path, and treat a passthrough and a judge failure in the same session as one incident,
since together they degrade both what the participant reads and how the session is evaluated.

\subsection{What to Store}\label{sec:disc-storage}

Storing only English keeps research data independent of language, which is why the boundary was
designed that way. It also discards evidence. When a participant's explanation is mistranslated,
the stored record cannot show whether the fault lay with the participant or with the translation,
and native-language evaluators have nothing to inspect. Studies that use this design should consider
storing the original language and text, the English translation and translation metadata alongside
each other, with audio still discarded, subject to participant consent.

\subsection{Implications for the Two Teach2Learn Systems}

For the fine-tuned portal, the five-model run (Appendix~\ref{app:models}) points to Qwen~2.5 7B,
which no subject currently uses, as the most stable learner. For the RAG tool, the priority is
repetition: its student prompt tells the learner to assume the explanation was helpful enough, which
works against the probing LbT relies on, and revising it is a larger lever than any change to the
translation layer.

\trackb{Relate the benchmark findings to what participants reported: whether teaching in their own
language changed their experience, and whether the Review-tier languages were acceptable in
practice.}

%% =======================================================================================
\section{Limitations}\label{sec:limitations}
%% =======================================================================================

\subsection{Limitations of System Design}\label{sec:lim-system}

\paragraph{One translation service.} All multilingual quality depends on one commercial service,
which returns Sindhi in the wrong script and translates the same sentence differently on different
calls, most often for low-resource languages.

\paragraph{Silent degradation.} By design, a translation failure in a live session passes the
original text through. The benchmark harness detects and retries such failures, but the live module
at the time of the benchmark did not retry, and its default timeout dropped roughly one crossing in
sixty under load in an earlier run. A participant whose answer passes
through untranslated sends non-English text to an English judge, which can degrade feedback and, in
the fine-tuned portal, lengthen the session.

\paragraph{Interface language.} The fine-tuned portal translates the conversation and feedback but
not its buttons, headings, survey questions or rating labels, so a participant teaching in Tamil
does so inside an English interface.

\paragraph{Grammatical gender.} Hindi and several other languages mark the addressee's gender on
the verb. The translation service chooses that form without knowing the participant's gender. In
the example of Figure~\ref{fig:example} it addresses the teacher in the feminine; we did not
measure how often this happens across languages or turns.

\paragraph{English-only storage.} The participant's original wording is not kept, so neither the
fidelity of the inbound translation nor the participant's own phrasing can be recovered
(Section~\ref{sec:disc-storage}).

\paragraph{Speech input is unevaluated.} We did not measure recognition accuracy in any language;
the coverage statements in Section~\ref{sec:voice} are the providers'. Transcription errors would
compound with translation errors before the learner sees the answer.

\paragraph{Corpus gap.} One of the five Computer Programming textbooks could not be parsed and is
absent from the RAG tool's retrieval index.

\subsection{Limitations of the Benchmark}\label{sec:lim-benchmark}

\paragraph{An architectural comparison.} The RAG learner is GPT-4o; the fine-tuned learners are
7--9B-parameter models, and the two also differ in prompt, context and self-evaluation. A RAG
advantage cannot be credited to retrieval, and a fine-tuned advantage cannot be credited to
fine-tuning.

\paragraph{Language effects are not isolated.} Every judged fluency score combines inbound
translation, the learner, outbound translation and the judge, so the agreement between systems is
consistent with, but does not prove, a dominant role for the translation layer.

\paragraph{An LLM judge.} Fidelity, fluency and depth are one model's labels. The judge agrees with
itself, but no native reader has confirmed any language's scores; LLM judges are known to score high
in low-resource and non-Latin-script languages \cite{hada2024multilingual}; and the fluency difference
we report is smaller than the judge's disagreement with itself.

\paragraph{Inbound fidelity is not measured.} Fidelity grades only the outward translation of the
teacher's sentence. The translation back into English, which is what the learner and judge
actually read, is checked for passthrough and script leaks but not for meaning. In the second turn
of the lesson in Figure~\ref{fig:example}, the teacher said that round-robin stops any one process
from starving the others. The English that came back said, in the fine-tuned run, that no process can
pre-empt another ``unless it has caught a CPU burst'' (its Hindi scored 5 for fidelity) and, in the
RAG run, ``while holding the CPU''. Neither distortion registered in any score.

\paragraph{Depth measures phrasing.} The largest difference in the RAG learner's favour, question
depth, is confounded with its uniform question template (Section~\ref{sec:res-depth}).

\paragraph{Scripted, machine-translated teachers.} Teacher input is three reviewed English sentences
per lesson, machine-translated into each language. Real participants write longer, less tidy text
that mixes English technical terms into their own language, and a machine-translated sentence is
likely easier to translate back than one a native speaker writes.

\paragraph{Decoding and scale.} The two sides used temperatures of 0.2 and 0.25. Sweep B used eight of
the 24 lessons with one repeat per language, and every lesson has three turns, fewer than a real
session.

\subsection{Limitations of User Study}\label{sec:lim-study}

\trackb{To be written by the study team.}

%% =======================================================================================
\section{Conclusion and Future Work}\label{sec:conclusion}
%% =======================================================================================

An English-only learning-by-teaching tutor can be made multilingual without rebuilding it. We drew a
single translation boundary around two independent Teach2Learn implementations, one with fine-tuned
small-model learners and one with a retrieval-augmented GPT-4o learner. Only what the participant
writes and reads is translated; inference, evaluation and stored research data remain English, and
dictated speech joins the same path as typed text. Across 1,575 boundary crossings per system, no text
passed untranslated and no Indian-language script reached the English side.

Sharing the boundary made the two architectures comparable on identical input. They agreed on the
quality tier of all 22 languages and their per-language fluency was closely correlated, suggesting
that the translation layer is a dominant determinant of multilingual quality: 15 languages are ready,
six need native review, and Sindhi cannot be supported with the current service. With language held
constant, the two learners failed in different ways. The retrieval-augmented learner wrote slightly
more fluent messages but repeated a question in a third of its conversations; the fine-tuned learners
followed the teacher more closely and never repeated one. Two pre-registered explanations for the
repeats were rejected. \trackb{One or two sentences on what the user study adds.}

\subsection{Future Work}

\begin{itemize}
\item Test generation-side explanations for the RAG learner's repeats by revising its student prompt
(removing the instruction to assume the explanation was sufficient, and discouraging reuse of earlier
questions), then re-running the benchmark.
\item Compare the translation pipeline with a natively multilingual learner on the same benchmark.
\item Measure transcription accuracy per language for both speech services.
\item Measure inbound translation fidelity, and add a second translation service as a control.
\item Translate the portal's interface, choose grammatical gender from the participant's preference,
and restore Sindhi once a service returns it in Perso-Arabic script.
\end{itemize}

\backmatter

\bmhead{Acknowledgements}
\todo{Acknowledgements.}

\section*{Declarations}

\begin{itemize}
\item Funding \todo{}
\item Conflict of interest/Competing interests \todo{}
\item Ethics approval and consent to participate \trackb{IRB reference for the user study.}
\item Consent for publication \todo{}
\item Data availability: The lesson bank, the benchmark results files for both systems, and the
question-depth labels are available from the corresponding author on reasonable request.
\todo{Replace with a public archive link if the repository is released.}
\item Materials availability \todo{}
\item Code availability: The translation and transcription modules, both benchmark harnesses and
the analysis and figure scripts are available from the corresponding author on reasonable request.
\todo{Replace with a public repository link if released.}
\item Use of AI tools: AI assistance in writing the software and in drafting this article is
described in Section~\ref{sec:ai-use}. No AI tool is listed as an author.
\item Author contribution \todo{}
\end{itemize}

\begin{appendices}

\section{Prompt Templates}\label{app:prompts}

\subsection{Judge: Fidelity}
Used identically for both systems. \texttt{\{language\}} is the language name.
\begin{lstlisting}
You compare a source English sentence with its translation into {language}.
Reply with exactly two lines and nothing else:
Score: <1-5>
Note: <one short clause>
5 means the translation preserves the full technical meaning in correct {language} script.
1 means the meaning is lost, the script is wrong, or the text is still English.
\end{lstlisting}
The user message is \texttt{English source:} followed by the teacher sentence, then
\texttt{\{language\} translation:} followed by the seed translation.

\subsection{Judge: Fluency (Aligned Version)}
\begin{lstlisting}
You judge one piece of text written in {language}, addressed by a student to their teacher.
Reply with exactly two lines and nothing else:
Score: <1-5>
Note: <one short clause>
5 means natural, fluent {language} in the correct script that a learner would actually write.
1 means unintelligible, wrong script, or untranslated English.
\end{lstlisting}
The user message is \texttt{\{language\} text:} followed by the learner's outbound message.

\subsection{Judge: Question Category}
The user message is \texttt{Question:} followed by the question sentence alone.
\begin{lstlisting}
You classify one question a student asked while learning a computer-science topic, using Graesser and Person's question categories.
Classify by the information the question seeks, not by its politeness form: "Could you explain how X works?" seeks a mechanism.

verification: whether something is true (yes/no)
disjunctive: which of the stated alternatives
concept completion: who, what, where or when, to fill in a missing detail
feature specification: the properties or attributes of something
quantification: how much or how many
definition: what a term or concept means
example: an example or instance of something
comparison: how two things are alike or differ
interpretation: what can be concluded from a pattern, data or situation
causal antecedent: what causes or caused something
causal consequence: what follows or would follow from something
goal orientation: why something is done, its purpose
instrumental/procedural: how something is done or works, its steps or mechanism
enablement: what makes something possible
expectational: why something expected does not happen
judgmental: an evaluation, opinion or recommendation
assertion: a statement of not knowing rather than a question
request/directive: asks the listener to do something unrelated to seeking information

Reply with exactly one line and nothing else:
Category: <one category name from the list>
\end{lstlisting}

\subsection{RAG Learner}
As run in the benchmark, with retrieved passages in \texttt{\{context\}}. Typographic quotes and
dashes are shown in ASCII. The teacher-history arm (Section~\ref{sec:prereg}) inserts one block
immediately before \texttt{Student Context:}; it is empty on the first turn, so that turn's prompt is
byte-identical to the one below. On later turns it reads:
\begin{lstlisting}
Earlier in this lesson, the teacher explained:
1. {first teacher sentence}
2. {second teacher sentence}
\end{lstlisting}
\begin{lstlisting}
You are a curious student learning from the teacher. You do not have any prior knowledge of the topic other than the Student Context and Course Material provided below.

Course Material (retrieved from the textbook, for your reference only - never tell the
teacher you have a textbook; you are learning only from what they explain):
{context}

Student Context:
{student_memory}

Teacher explanation:
"{teacher_explanation}"

Your task:
- Decide how well you understood the explanation.
  * If it made sense overall, even if it wasn't detailed, -> rating = "understood", message = null
  * If you're still very confused -> rating = "confused", but ask ONE short and polite clarifying question
  * If the Course Material above contradicts or is clearly more complete than what the teacher said, treat that as still-confused and gently ask about the specific gap, without ever citing or naming the material itself.
- Write a short reflection about your understanding in natural language, e.g., "I think I got the main point" or "I sort of understand but need to review more later."
- If you noticed anything missing, keep it very light and optional (e.g., "maybe more detail later could help"), but only if really needed.

**Always show appreciation for the teacher's effort. Assume the explanation was helpful enough for now, even if not perfect.**

Respond ONLY in valid JSON that matches this schema:
{
  "message": "Your follow-up question or null",
  "rating": "understood|confused",
  "reflection": "How you understood the concept",
  "missing_points": ["point1", "point2"]
}
\end{lstlisting}

\subsection{Fine-Tuned Learner Input Format}
The adapters take no free-text instruction at inference beyond the format they were trained on. The
first turn is a system trigger; later teacher turns are rendered as:
\begin{lstlisting}
Previous Question: {learner's previous question}
Teacher's Answer: {teacher's English answer}
\end{lstlisting}

\section{Lesson Bank}\label{app:lessons}

Version 1, reviewed 2026-09-08. $^{*}$Carried over verbatim from an earlier
single-lesson-per-subject pilot, so that cells can be compared with it. One of these, on heaps,
describes the root as the smallest element, which holds for a min-heap only, and uses the bank's only
raw arithmetic (\texttt{2i+1}, \texttt{2i+2}); it was kept unchanged to preserve that comparison.
$^{\ddagger}$Used in Sweep B. All 24 lessons are used in Sweep A.

\footnotesize
\begin{longtable}{p{0.3\textwidth}p{0.64\textwidth}}
\toprule
Lesson & Teacher script \\
\midrule
\endhead
\texttt{cp/c-basics/types-variables-constants} & (1) Every variable in C has a type fixed at compile time, because the compiler has to know how many bytes to reserve and how to interpret the bits it finds there. (2) An int and a float of the same width store numbers in completely different layouts, so assigning one to the other converts the value rather than reusing the bits. (3) Declaring a value const tells the compiler it will never be assigned again, which lets it reject a whole class of mistakes before the program is ever run. \\ \midrule
\texttt{cp/c-basics/operators-and-expressions} & (1) Operator precedence decides how an expression without brackets is grouped, so a star binds more tightly than a plus and the multiplication happens first. (2) Integer division in C throws away the remainder instead of rounding, which is why seven divided by two is three and not three point five. (3) The logical and operator stops evaluating as soon as its left side is false, so the right side can safely depend on the check that came before it. \\ \midrule
\texttt{cp/control-flow/conditions-and-switch} & (1) An if statement picks one branch by evaluating a condition, and in C any non-zero value counts as true rather than only a dedicated boolean. (2) A switch compares one value against several constant labels, which reads more clearly than a long chain of else-if tests on the same variable. (3) Each case falls through into the next unless it ends with a break, so a missing break silently runs the code for every label below it. \\ \midrule
\texttt{cp/control-flow/loops-and-sentinel-values}$^{\ddagger}$ & (1) A for loop is the right choice when the number of repetitions is known in advance, because the counter, the test, and the update all sit in one header. (2) A while loop suits the case where the program keeps reading until some condition changes, since the number of iterations is not known before it starts. (3) A sentinel is a reserved value that marks the end of the input, so the loop can stop on the data itself instead of being told the count beforehand. \\ \midrule
\texttt{cp/pointers/addresses-and-pointers} & (1) A pointer is a variable whose value is the address of another variable, so it tells you where something lives rather than what it is. (2) The ampersand operator takes the address of a variable and the star operator goes back the other way to read the value stored at an address. (3) Passing an address into a function lets that function change the caller's variable, because C otherwise copies every argument and the original stays untouched. \\ \midrule
\texttt{cp/pointers/dynamic-memory-allocation}$^{*}$$^{\ddagger}$ & (1) malloc reserves a block of memory on the heap while the program is running and returns a pointer to its first byte, because the size is not known at compile time. (2) Every successful malloc must be matched by exactly one free, otherwise the block stays reserved for the whole life of the program and the process leaks memory. (3) Using a pointer after its block has been freed is undefined behaviour, so it is good practice to set the pointer to NULL immediately after freeing it. \\ \midrule
\texttt{dsa/analysis/asymptotic-notation}$^{\ddagger}$ & (1) Big-O describes how the running time grows as the input grows, which lets two algorithms be compared without measuring them on any particular machine. (2) Constant factors and lower-order terms are dropped, so three n squared plus five n is simply order n squared once the input is large enough. (3) Big-O is an upper bound rather than an exact cost, which is why an algorithm that is order n squared in the worst case can still be fast on typical inputs. \\ \midrule
\texttt{dsa/analysis/recurrences} & (1) A recurrence expresses the cost of a recursive algorithm in terms of the cost of the smaller calls it makes, plus the work done to combine them. (2) Merge sort gives the recurrence T of n equals two T of n over two plus n, because it splits the array in half and then merges in linear time. (3) Solving that recurrence gives n log n, since the splitting produces log n levels and every level does a total of n work. \\ \midrule
\texttt{dsa/linear/stacks} & (1) A stack only allows access at one end, so the last item pushed is always the first item popped and nothing in the middle can be reached directly. (2) That restriction is what makes it the natural structure for undo history and for the call stack, where the most recent call must finish first. (3) Both push and pop touch only the top element, so each costs constant time no matter how many items the stack is holding. \\ \midrule
\texttt{dsa/linear/queues-and-deques} & (1) A queue adds at one end and removes at the other, so items leave in the order they arrived and no item can overtake the ones ahead of it. (2) Implementing a queue as a circular array avoids shifting every element after each removal, because the front and back indices wrap around instead. (3) A deque relaxes the rule and allows insertion and removal at both ends, which makes it a superset of both the stack and the queue. \\ \midrule
\texttt{dsa/trees-heaps/tree-terminology-and-traversals} & (1) A tree is a set of nodes where every node except the root has exactly one parent, so there is precisely one path from the root to any node. (2) The height of a tree is the length of the longest root-to-leaf path, and it bounds the cost of any operation that walks down from the root. (3) An in-order traversal of a binary search tree visits the keys in sorted order, because it fully explores the left subtree before touching the node itself. \\ \midrule
\texttt{dsa/trees-heaps/heaps}$^{*}$$^{\ddagger}$ & (1) A binary heap is a complete binary tree where every parent is ordered against its children, so the smallest element is always sitting at the root. (2) Because the tree is complete, the heap can be stored in a plain array, where the children of index i live at 2i+1 and 2i+2 and no pointers are needed. (3) Inserting or removing an element only walks one path between a leaf and the root, so both operations cost logarithmic time in the number of elements. \\ \midrule
\texttt{oop\_java/classes/constructors} & (1) A constructor runs once when an object is created, and its job is to leave the new object in a state where every method can safely be called. (2) Java supplies a no-argument constructor only when the class declares none, so writing any constructor at all removes that default. (3) Overloaded constructors let the same class be built from different inputs, and one can call another with this so the shared setup is written only once. \\ \midrule
\texttt{oop\_java/classes/access-control-and-object-identity} & (1) Declaring a field private stops code outside the class from touching it directly, so the class keeps control over every change to its own state. (2) The double-equals operator compares references rather than contents, so two distinct objects holding identical values are still not equal by that test. (3) Overriding equals means overriding hashCode as well, because collections like HashMap decide which bucket to search using the hash before they compare at all. \\ \midrule
\texttt{oop\_java/oop-design/inheritance} & (1) Inheritance lets a subclass reuse the fields and methods of its superclass, so shared behaviour is written once and specialised only where it differs. (2) A subclass constructor always runs its superclass constructor first, because the inherited part of the object has to be initialised before the new part. (3) Java allows only one superclass per class, which avoids the ambiguity that arises when the same method is inherited along two different paths. \\ \midrule
\texttt{oop\_java/oop-design/polymorphism}$^{\ddagger}$ & (1) A superclass reference can point at any subclass object, so one variable can hold different concrete types over the life of the program. (2) Which overriding method runs is decided at run time from the object's actual class, not at compile time from the type of the reference. (3) That is what lets a single loop over a list of shapes call draw on each one without knowing or caring which shape it is holding. \\ \midrule
\texttt{oop\_java/concurrency/threads} & (1) A thread is a separate flow of execution inside the same process, so all threads share the heap while each keeps its own stack. (2) Calling start hands the work to the scheduler and returns immediately, whereas calling run directly just executes the body on the current thread. (3) The order in which two independent threads make progress is not guaranteed, so a program that only works under one interleaving is already broken. \\ \midrule
\texttt{oop\_java/concurrency/synchronization-and-locks}$^{*}$$^{\ddagger}$ & (1) When two threads update the same field at once one update can be lost, because reading, changing, and writing back is not a single indivisible step. (2) A synchronized block lets only one thread hold the object's monitor at a time, so every other thread waits until that block finishes. (3) Acquiring two locks in different orders in different threads can deadlock, so the usual rule is to always take locks in one fixed global order. \\ \midrule
\texttt{os/processes/process-lifecycle} & (1) A process moves between the ready, running, and blocked states, and only the running state is actually holding a CPU at that moment. (2) A process that asks for input goes from running to blocked rather than staying on the CPU, because otherwise it would waste the processor doing nothing. (3) The operating system keeps a control block for each process holding its registers and program counter, which is what makes it possible to resume it later. \\ \midrule
\texttt{os/processes/interprocess-communication} & (1) Separate processes have separate address spaces, so one cannot simply read the other's variables and the operating system has to provide a channel. (2) A pipe carries a stream of bytes from one process to another, and the reader blocks when the pipe is empty rather than seeing partial data. (3) Message passing is easier to reason about than shared memory because the kernel handles the synchronisation, though it costs a copy on every exchange. \\ \midrule
\texttt{os/threads-scheduling/multithreading} & (1) Threads within one process share the same address space, so switching between them is cheaper than switching between two separate processes. (2) That sharing is also the danger, because two threads writing the same variable can interleave in ways neither of them was written to expect. (3) A program that spends most of its time waiting on input and output can gain from threads even on a single core, since one thread runs while another waits. \\ \midrule
\texttt{os/threads-scheduling/cpu-scheduling}$^{*}$$^{\ddagger}$ & (1) The scheduler decides which ready process runs next on the CPU, because there are usually more runnable processes than cores. (2) Round-robin gives each process a fixed time quantum in turn, so no single process can starve the others while it holds the CPU. (3) A very small quantum wastes time on context switches, while a very large one makes the scheduler behave like plain first-come first-served. \\ \midrule
\texttt{os/virtual-memory/demand-paging}$^{\ddagger}$ & (1) Demand paging loads a page into memory only when it is first touched, so a program can start running before all of it has been read from disk. (2) Touching a page that is not resident raises a page fault, and the operating system fetches it and restarts the instruction that faulted. (3) This lets the total size of running programs exceed physical memory, because only the pages actually being used need to be resident at any moment. \\ \midrule
\texttt{os/virtual-memory/page-replacement} & (1) When memory is full and a new page is needed, the operating system has to choose an existing page to evict to make room for it. (2) Least-recently-used approximates the ideal choice by assuming a page untouched for a long time is unlikely to be needed soon. (3) Evicting a page that has been modified costs an extra write back to disk, which is why a clean page is usually the cheaper victim. \\
\botrule
\end{longtable}
\normalsize

\section{Results by Language and Subject}\label{app:matrices}

Each cell is fidelity/fluency, averaged over that language's judged turns in that subject. Subject
also selects the fine-tuned model. $^\dagger$One lesson never produced data (the Punjabi schema
payload). $^{*}$The service returns Devanagari; Perso-Arabic is the dominant script for Sindhi.

\begin{table}[p]
\caption{Fine-tuned learners.}\label{tab:ft-matrix}
\footnotesize
\begin{tabular*}{\textwidth}{@{\extracolsep\fill}lllllllll}
\toprule
Language & Script & Code & \texttt{cp} & \texttt{dsa} & \texttt{oop\_java} & \texttt{os} & Mean & Tier \\
\midrule
Telugu & Telugu & \texttt{te} & 4.5/4.8 & 4.3/4.8 & 4.7/4.7 & 4.5/4.8 & 4.5/4.8 & Ready \\
Bengali & Bengali & \texttt{bn} & 4.2/5.0 & 4.3/4.8 & 4.8/4.8 & 4.2/4.5 & 4.4/4.8 & Ready \\
Punjabi$^\dagger$ & Gurmukhi & \texttt{pa} & 5.0/5.0 & 4.7/5.0 & 4.7/4.5 & 4.5/4.7 & 4.7/4.8 & Ready \\
Gujarati & Gujarati & \texttt{gu} & 4.8/5.0 & 4.7/4.8 & 4.5/4.3 & 4.8/4.8 & 4.7/4.8 & Ready \\
Hindi & Devanagari & \texttt{hi} & 5.0/4.5 & 4.2/4.7 & 4.8/4.7 & 4.3/4.7 & 4.6/4.6 & Ready \\
Tamil & Tamil & \texttt{ta} & 4.2/4.8 & 4.8/4.8 & 4.7/4.3 & 4.3/4.3 & 4.5/4.6 & Ready \\
Malayalam & Malayalam & \texttt{ml} & 5.0/4.7 & 4.7/4.5 & 4.5/4.3 & 4.8/4.8 & 4.8/4.6 & Ready \\
Assamese & Bengali & \texttt{as} & 4.5/4.7 & 3.8/4.8 & 4.3/4.7 & 4.2/4.2 & 4.2/4.6 & Ready \\
Nepali & Devanagari & \texttt{ne} & 3.8/5.0 & 4.7/4.5 & 4.7/4.2 & 4.5/4.7 & 4.4/4.6 & Ready \\
Marathi & Devanagari & \texttt{mr} & 4.3/4.5 & 4.5/4.5 & 4.5/4.7 & 4.5/4.5 & 4.5/4.5 & Ready \\
Urdu & Perso-Arabic & \texttt{ur} & 5.0/4.7 & 4.5/4.3 & 4.7/4.7 & 4.5/4.5 & 4.7/4.5 & Ready \\
Kannada & Kannada & \texttt{kn} & 4.7/5.0 & 4.3/4.3 & 4.7/4.5 & 4.3/4.2 & 4.5/4.5 & Ready \\
Odia & Odia & \texttt{or} & 4.8/4.7 & 4.7/4.3 & 4.8/4.0 & 4.7/4.3 & 4.8/4.3 & Ready \\
Maithili & Devanagari & \texttt{mai} & 4.7/4.7 & 4.7/4.3 & 4.7/4.0 & 4.5/4.3 & 4.6/4.3 & Ready \\
Dogri & Devanagari & \texttt{doi} & 4.2/4.2 & 4.8/3.7 & 4.5/4.0 & 4.0/4.2 & 4.4/4.0 & Ready \\
Bodo & Devanagari & \texttt{brx} & 4.8/4.0 & 5.0/3.8 & 4.2/3.3 & 5.0/4.0 & 4.8/3.8 & Review \\
Kashmiri & Perso-Arabic & \texttt{ks} & 4.5/4.2 & 4.3/3.5 & 3.3/3.8 & 4.3/3.7 & 4.1/3.8 & Review \\
Manipuri & Meetei Mayek & \texttt{mni} & 5.0/4.0 & 4.5/3.7 & 4.8/3.3 & 4.8/3.7 & 4.8/3.7 & Review \\
Santali & Ol Chiki & \texttt{sat} & 4.5/3.7 & 4.7/3.0 & 4.8/3.7 & 5.0/3.3 & 4.8/3.4 & Review \\
Konkani & Devanagari & \texttt{kok} & 4.5/3.3 & 4.8/3.5 & 4.7/3.3 & 4.3/3.3 & 4.6/3.4 & Review \\
Sanskrit & Devanagari & \texttt{sa} & 3.3/3.2 & 3.2/3.2 & 3.2/2.7 & 3.0/3.0 & 3.2/3.0 & Review \\
Sindhi & Devanagari$^{*}$ & \texttt{sd} & 1.3/1.7 & 1.5/2.2 & 1.0/2.3 & 1.5/2.3 & 1.3/2.1 & Blocked \\
\botrule
\end{tabular*}
\end{table}

\begin{table}[p]
\caption{RAG learner.}\label{tab:rag-matrix}
\footnotesize
\begin{tabular*}{\textwidth}{@{\extracolsep\fill}lllllllll}
\toprule
Language & Script & Code & \texttt{cp} & \texttt{dsa} & \texttt{oop\_java} & \texttt{os} & Mean & Tier \\
\midrule
Malayalam & Malayalam & \texttt{ml} & 5.0/4.8 & 4.8/4.8 & 4.5/5.0 & 5.0/5.0 & 4.8/4.9 & Ready \\
Punjabi$^\dagger$ & Gurmukhi & \texttt{pa} & 4.3/5.0 & 4.8/5.0 & 4.7/4.8 & 4.3/5.0 & 4.6/5.0 & Ready \\
Gujarati & Gujarati & \texttt{gu} & 4.8/4.8 & 4.7/4.8 & 4.5/4.7 & 4.8/4.8 & 4.7/4.8 & Ready \\
Urdu & Perso-Arabic & \texttt{ur} & 4.8/5.0 & 4.7/4.8 & 4.7/4.8 & 4.5/4.7 & 4.7/4.8 & Ready \\
Hindi & Devanagari & \texttt{hi} & 5.0/5.0 & 4.3/4.7 & 4.7/4.5 & 4.5/5.0 & 4.6/4.8 & Ready \\
Telugu & Telugu & \texttt{te} & 4.2/4.8 & 4.0/4.7 & 4.7/4.8 & 4.5/5.0 & 4.3/4.8 & Ready \\
Odia & Odia & \texttt{or} & 4.5/4.7 & 4.5/4.7 & 4.7/4.5 & 4.5/5.0 & 4.5/4.7 & Ready \\
Nepali & Devanagari & \texttt{ne} & 4.0/5.0 & 4.7/4.7 & 4.7/4.5 & 4.3/4.7 & 4.4/4.7 & Ready \\
Kannada & Kannada & \texttt{kn} & 4.7/4.5 & 4.3/4.7 & 4.7/4.8 & 4.5/4.7 & 4.5/4.7 & Ready \\
Bengali & Bengali & \texttt{bn} & 4.3/4.5 & 4.3/4.8 & 4.7/4.8 & 4.2/4.7 & 4.4/4.7 & Ready \\
Marathi & Devanagari & \texttt{mr} & 4.3/4.7 & 4.3/4.7 & 4.5/4.7 & 4.3/4.7 & 4.4/4.7 & Ready \\
Assamese & Bengali & \texttt{as} & 4.7/4.5 & 4.7/4.5 & 4.5/4.8 & 4.3/4.7 & 4.5/4.6 & Ready \\
Tamil & Tamil & \texttt{ta} & 4.0/5.0 & 4.5/4.7 & 4.7/4.5 & 4.3/4.2 & 4.4/4.6 & Ready \\
Maithili & Devanagari & \texttt{mai} & 4.7/4.5 & 4.7/4.5 & 4.8/4.2 & 4.2/4.3 & 4.6/4.4 & Ready \\
Dogri & Devanagari & \texttt{doi} & 4.2/4.2 & 4.7/4.5 & 4.5/4.2 & 4.3/4.2 & 4.4/4.2 & Ready \\
Bodo & Devanagari & \texttt{brx} & 4.8/4.0 & 4.8/4.0 & 4.3/3.8 & 5.0/4.0 & 4.8/4.0 & Review \\
Kashmiri & Perso-Arabic & \texttt{ks} & 4.7/4.0 & 4.7/4.2 & 4.3/3.5 & 4.3/4.2 & 4.5/4.0 & Review \\
Konkani & Devanagari & \texttt{kok} & 4.3/3.7 & 4.8/3.8 & 4.7/3.8 & 4.5/4.0 & 4.6/3.8 & Review \\
Manipuri & Meetei Mayek & \texttt{mni} & 5.0/4.0 & 4.5/3.7 & 5.0/3.7 & 5.0/3.8 & 4.9/3.8 & Review \\
Santali & Ol Chiki & \texttt{sat} & 4.8/3.3 & 4.8/3.6 & 4.8/3.3 & 5.0/4.0 & 4.9/3.6 & Review \\
Sanskrit & Devanagari & \texttt{sa} & 3.5/3.2 & 3.3/3.2 & 3.2/3.5 & 3.0/3.2 & 3.2/3.2 & Review \\
Sindhi & Devanagari$^{*}$ & \texttt{sd} & 1.0/2.0 & 1.0/1.8 & 1.0/1.7 & 1.0/2.0 & 1.0/1.9 & Blocked \\
\botrule
\end{tabular*}
\end{table}

Figure~\ref{fig:heatmap} shows the fluency matrices as a heatmap.

\begin{figure}[h]
\centering
\includegraphics[width=0.92\textwidth]{figures/fig3-heatmap-fluency.pdf}
\caption{Judged fluency by language and subject for both systems (1--5; darker is more fluent).
Rows are ordered by mean fluency across both systems. Subject also selects the fine-tuned
model.}\label{fig:heatmap}
\end{figure}

\section{Composite Score and Weighting Sensitivity}\label{app:composite}

The pre-registered analysis combined the six dimensions of Table~\ref{tab:dimensions} into a weighted
0--100 composite. We report it here rather than in the main text because its weighting is necessarily
normative: the verdict depends on how much each dimension is judged to matter. With the pre-registered
weights (fluency 36, question depth 20, lexical variety 12, non-repetition 12, robustness 12,
reliability 8) the RAG learner scores 81.4 [78.8, 83.9] and the fine-tuned learners 75.0 [72.6, 77.3],
a gap of $+6.5$ [$+3.6$, $+9.6$]. That gap rests on question depth: without it, the five remaining
dimensions give 79.3 against 78.6, a gap of $+0.7$ [$-2.2$, $+3.3$], which is a tie.
Table~\ref{tab:weights} recomputes the composite under other weightings, including both superseded
versions. RAG leads under every weighting that includes depth except one that prioritises
conversation over translation; without depth the systems tie. The first, pre-declared version had
a single variety dimension of weight 30, fixed before any transcripts were read; it rewarded a learner
for varying the first six words of questions it then repeated verbatim, so it was split evenly into
lexical variety and non-repetition. Question depth was then added under its pre-registration.

\begin{table}[h]
\caption{Composite under alternative weightings (fluency, depth, lexical variety, non-repetition,
robustness, reliability).}\label{tab:weights}
\footnotesize
\begin{tabular*}{\textwidth}{@{\extracolsep\fill}llrrll}
\toprule
Weighting & Weights & FT & RAG & RAG$-$FT [95\% CI] & Leader \\
\midrule
This paper & 36/20/12/12/12/8 & 75.0 & 81.4 & $+6.5$ [$+3.6$, $+9.6$] & RAG \\
Equal & 16/20/16/16/16/16 & 76.0 & 81.5 & $+5.5$ [$+2.1$, $+9.1$] & RAG \\
Multilingual first & 48/20/8/8/12/4 & 74.3 & 81.5 & $+7.2$ [$+4.6$, $+10.0$] & RAG \\
Conversation first & 20/20/8/32/12/8 & 80.2 & 78.7 & $-1.5$ [$-6.6$, $+3.4$] & Tie \\
Without depth (revision 1) & 45/---/15/15/15/10 & 78.6 & 79.3 & $+0.7$ [$-2.2$, $+3.3$] & Tie \\
Pre-declared (superseded) & 45/---/30$^{a}$/---/15/10 & 71.6 & 80.9 & $+9.3$ [$+5.9$, $+12.7$] & RAG \\
\botrule
\end{tabular*}
\footnotesize{$^{a}$Single ``variety'' dimension measured by opener entropy alone.}
\end{table}

\section{Model Run on the Fine-Tuned Side}\label{app:models}

Running all five deployed models on the same \texttt{os} and \texttt{cp} lessons separates model from
subject. Qwen~2.5 7B was the only model stable across both subjects: the share of follow-ups beginning
with its most common opening was 33\% and 30\%, its mean Bloom level 1.00 and 0.98, and it reached the
Apply level 17 and 14 times against 0--9 for the others, while also generating fastest (2.7\,s per
generation). Qwen3 8B collapsed to four distinct openings on \texttt{cp}, beginning 85\% of follow-ups
with the same phrase, against 22\% on \texttt{os}; the \texttt{cp} figure reproduced in an independent
run a day apart (83\%). The \texttt{cp} lessons reduced opener variety for every model, so Qwen3 is
the extreme of a general effect. Granite 3.3 8B on \texttt{os}, run twice a day apart, gave 9\% both
times, which bounds run-to-run noise. No subject currently routes to Qwen~2.5.

\section{Teacher-History Experiment: All Measures}\label{app:experiment}

Table~\ref{tab:experiment} compares the committed RAG run with the teacher-history arm, paired by lesson. No measure moved.

\begin{table}[h]
\caption{Teacher-history experiment, paired by lesson (committed RAG run vs.\ the new arm).}\label{tab:experiment}
\footnotesize
\begin{tabular*}{\textwidth}{@{\extracolsep\fill}lrrrl}
\toprule
Measure & RAG & + history & Change & 95\% CI \\
\midrule
Conversations repeating a question (primary) & 23/72 & 24/72 & $+1.4$ pts & [$-10.8$, $+14.2$] \\
Looks back to an earlier sentence & 30\% & 31\% & $+1.3$ pts & [$-9.7$, $+12.7$] \\
Answers the current sentence & 44\% & 41\% & $-3.4$ pts & [$-12.2$, $+4.4$] \\
Follow-ups that are questions & 78\% & 82\% & $+3.7$ pts & [$-2.7$, $+11.2$] \\
Questions rated deep & 90\% & 92\% & $+1.6$ pts & [$-2.0$, $+5.5$] \\
Lexical variety (entropy, rarefied to 169) & 0.78 & 0.80 & $+0.011$ & [$-0.028$, $+0.056$] \\
\botrule
\end{tabular*}
\end{table}

\section{Differences Recorded but Not Scored}\label{app:notscored}

These differences follow from how each system was built or prompted; scoring them would compare design choices rather than learners.

\begin{table}[h]
\caption{Differences recorded but not scored.}\label{tab:notscored}
\footnotesize
\begin{tabular*}{\textwidth}{@{\extracolsep\fill}p{0.34\textwidth}rrp{0.36\textwidth}}
\toprule
Observation & FT & RAG & Why not scored \\
\midrule
Follow-ups that are questions (English) & 100\% & 78\% & RAG prompt permits a reflection; adapter is trained to ask \\
Language turns answered with a reflection & 0/525 & 21/525 & Same prompt difference \\
Words per follow-up & 20.3 & 16.9 & Set by each system's instructions \\
Distinct follow-up strings & 99\% & 78\% & A fixed lead-in with a new ending counts as distinct \\
Runs flagged by the harness's opener check & 36/72 & 53/72 & RAG's reflections all begin alike \\
\botrule
\end{tabular*}
\end{table}

\section{Pre-Registration Text}\label{app:prereg}

All three entries are reproduced from the project's pre-registration file, written before the data
they judge existed. Only file paths and script names have been removed.

\paragraph{RAG teacher-history experiment. Written 2026-09-14T18:31:32Z.}
\emph{Change under test.} The RAG student receives the teacher's earlier sentences in the current
lesson, in addition to its own earlier replies. The live app already records these but never passes
them to the student chain. Nothing else changes: same prompt text plus one history block, same GPT-4o
via OpenRouter, temperature 0.25, top-$p$ 0.9, same $k=4$ retrieval, same 24 lessons $\times$ 3
repeats, and the same opening question per lesson, reused from the committed English run so openings
cannot differ between arms.
\emph{Primary metric (decided now).} English conversations containing a question asked twice, word
for word. Baseline RAG: 23 of 72. Fine-tuned: 0 of 72.
\emph{Prediction.} Falls from 23/72 toward 0/72. The prediction fails if the new arm's 95\% interval
for the change includes zero.
\emph{Secondary, reported whatever they show.} Looks-backward share, answers-current-sentence share,
lexical variety (rarefied opener entropy), share of follow-ups that are questions, question depth
(blinded judge, run after this).
\emph{Analysis.} Same scoring script, same seed 20260914. The new arm is added as a third column; the
committed RAG numbers are not replaced. No language sweep is run, so fluency for the new arm is not
measured, and any composite for it is labelled a projection that assumes RAG's fluency.

\paragraph{Question depth. Written 2026-09-14T18:35:26Z, before any depth score existed.}
\emph{Instrument.} Each follow-up question's question sentence (a stock lead-in such as ``That makes
sense, but \ldots'' is stripped) goes to the same judge model, \texttt{openai/gpt-5.6-luna}, alone:
no system name, no subject, no teacher sentence, and all three arms shuffled together. The judge names
one of the 18 question categories of Graesser and Person (1994). Grouping into depth levels is fixed
here: shallow (verification, disjunctive, concept completion, feature specification, quantification);
intermediate (definition, example, comparison, interpretation); deep (causal antecedent, causal
consequence, goal orientation, instrumental/procedural, enablement, expectational, judgmental); not a
question (assertion, request/directive), excluded from the denominator. Reflections (RAG replies with
no question) are excluded, as they are from lexical variety.
\emph{Reliability check.} A random 60 of the judged questions are judged a second time. Agreement on
depth level is reported next to the result. If it falls below 70\%, depth is reported but not scored.
\emph{Scored dimension.} ``Question depth'' = share of a system's questions classified deep. Weight
20. The five existing dimensions keep their relative weights, scaled by 0.8: fluency 36, lexical
variety 12, non-repetition 12, robustness 12, reliability 8. Verdict rule unchanged: the 95\% interval
of the difference must exclude zero. The previous five-dimension scorecard stays on the page.

\paragraph{Retrieval overlap and repeated questions. Written 2026-09-19T05:32:53Z, before any
retrieval was logged.} \emph{Question.} The pre-registered teacher-history experiment failed: RAG
repeated a question in 23 of 72 English conversations before the change and 24 of 72 after. The
remaining hypothesis is that retrieval causes the repeats. The student prompt tells the model to ask
about any gap between the teacher and the retrieved passages, so if two turns of a lesson retrieve the
same passages, the later turn may ask the same gap question again.
\emph{Data.} No new questions are generated. For each of the 24 lessons, each of its three teacher
sentences is passed to the retrieval step ($k=4$) against that subject's committed vectorstore,
exactly as the English sweep did, using the same embedding model. A passage is identified by its
source file, page and a hash of its text. Questions come from the committed RAG and RAG + teacher
history English runs. Retrieval depends only on the teacher sentence, so all repeats of a lesson, in
both arms, share its passages.
\emph{Unit, exposure, outcome.} Unit: a lesson and a pair of its turns, (1,2), (1,3), (2,3), giving
$24 \times 3 = 72$ units. Exposure: Jaccard overlap between the two turns' sets of four passage IDs.
Outcome: of the 6 conversations for that lesson (3 per arm), the share in which the later turn's
question sentence equals the earlier turn's, after the same normalisation the scorecard uses.
Conversations where either turn has no question sentence count as not repeating.
\emph{Primary test.} Repeat rate among units with overlap $> 0$, minus repeat rate among units with
overlap $= 0$. 95\% interval from a 5,000-draw bootstrap that resamples lessons (all three pairs of a
lesson move together), seed 20260914. If either group has fewer than 5 units, the primary test is
instead the Spearman correlation between overlap and repeat rate across the 72 units, with the same
lesson bootstrap.
\emph{Prediction and failure rule.} Prediction: units whose turns share retrieved passages repeat more
often. The hypothesis fails if the 95\% interval includes zero.
\emph{Reported whatever they show.} Spearman correlation between overlap and repeat rate; each arm
separately; the share of units with identical passage sets; the distribution of overlap values.

\section{Implementation Details}\label{app:implementation}

\paragraph{Translation module.} Both systems import one module, \texttt{translation.py}, verified
byte-identical with \texttt{cmp} before the comparison runs. It calls \texttt{sarvam-translate:v1},
which accepts 2,000 input characters, matching the portal's 2,000-character answer limit; an
alternative model with a 1,000-character limit would have dropped long answers back to untranslated
text. Longer text is split on sentence boundaries, and results are cached in process by (text,
source, target). Language codes are an explicit table rather than derived, because the service spells
Odia \texttt{od-IN}, not the ISO \texttt{or-IN}. A unit test locks the language picker and the code
table to the same key set, because a language present in the picker but missing a code would
otherwise fail as a silent passthrough. A pattern check rejects responses that contain fragments of
the service's API schema. The live module's default timeout was 5 seconds and it did not retry; the
benchmark harness made up to three attempts per crossing with backoff.

\paragraph{Interface.} In the RAG tool, which renders through Chainlit, the language is picked with
the same action-button mechanism as the subject, and interface messages the tool sends are translated.
The fine-tuned portal translates conversation and feedback but not its buttons, headings or survey.

\paragraph{Speech input.} The portal adds a microphone button that records with the browser's
\texttt{MediaRecorder}, re-encodes the clip to WAV in the browser, and uploads it to a session-scoped
endpoint with a 15\,MB limit; it places the returned text in the answer box for the participant to
review. The RAG tool uses Chainlit's built-in audio capture, wraps the streamed PCM chunks into a
24\,kHz WAV file on the server, and submits the text as the participant's turn. A shared
\texttt{transcription.py} module forwards the audio to Sarvam's Saaras v3 (default) or to Whisper
large-v3 via OpenRouter and discards it. Unlike translation, a transcription failure surfaces as an
error rather than passing anything through. The WAV re-encode exists because the default service
rejected the \texttt{audio/webm;codecs=opus} container that Chrome records by default. Whisper's
published language list omits Bodo, Dogri, Kashmiri, Konkani, Maithili, Manipuri, Odia and Santali.

\paragraph{Retrieval corpus.} Textbooks named in the course handouts were split into 1,000-character
chunks with 200-character overlap, embedded with \texttt{text-embedding-3-small}, and stored per
subject in Chroma (Table~\ref{tab:corpus}).

\begin{table}[h]
\caption{Retrieval corpus of the RAG tool. One of five \texttt{cp} sources (Forouzan and Gilberg)
failed to parse and is absent.}\label{tab:corpus}
\footnotesize
\begin{tabular*}{\textwidth}{@{\extracolsep\fill}llrr}
\toprule
Subject & Sources & Pages & Chunks \\
\midrule
\texttt{cp} & Hanly and Koffman, plus three reference books & 2,289 & 7,530 \\
\texttt{dsa} & Cormen et al.; Goodrich, Tamassia and Mount & 2,033 & 5,605 \\
\texttt{oop\_java} & Schildt; Horstmann & 1,788 & 5,148 \\
\texttt{os} & Silberschatz and Galvin & 983 & 3,281 \\
\midrule
Total & & 7,093 & 21,564 \\
\botrule
\end{tabular*}
\end{table}

\section{Supported Languages}\label{app:languages}

\begin{table}[h]
\caption{Teaching languages, scripts as returned by the translation service, and service codes.}\label{tab:language-codes}
\footnotesize
\begin{tabular*}{\textwidth}{@{\extracolsep\fill}llll}
\toprule
Language & Script & Code & Service code \\
\midrule
Hindi & Devanagari & \texttt{hi} & \texttt{hi-IN} \\
Tamil & Tamil & \texttt{ta} & \texttt{ta-IN} \\
Telugu & Telugu & \texttt{te} & \texttt{te-IN} \\
Kannada & Kannada & \texttt{kn} & \texttt{kn-IN} \\
Bengali & Bengali & \texttt{bn} & \texttt{bn-IN} \\
Malayalam & Malayalam & \texttt{ml} & \texttt{ml-IN} \\
Marathi & Devanagari & \texttt{mr} & \texttt{mr-IN} \\
Punjabi & Gurmukhi & \texttt{pa} & \texttt{pa-IN} \\
Gujarati & Gujarati & \texttt{gu} & \texttt{gu-IN} \\
Odia & Odia & \texttt{or} & \texttt{od-IN} (not \texttt{or-IN}) \\
Urdu & Perso-Arabic & \texttt{ur} & \texttt{ur-IN} \\
Assamese & Bengali & \texttt{as} & \texttt{as-IN} \\
Bodo & Devanagari & \texttt{brx} & \texttt{brx-IN} \\
Dogri & Devanagari & \texttt{doi} & \texttt{doi-IN} \\
Konkani & Devanagari & \texttt{kok} & \texttt{kok-IN} \\
Kashmiri & Perso-Arabic & \texttt{ks} & \texttt{ks-IN} \\
Maithili & Devanagari & \texttt{mai} & \texttt{mai-IN} \\
Manipuri & Meetei Mayek & \texttt{mni} & \texttt{mni-IN} \\
Nepali & Devanagari & \texttt{ne} & \texttt{ne-IN} \\
Sanskrit & Devanagari & \texttt{sa} & \texttt{sa-IN} \\
Santali & Ol Chiki & \texttt{sat} & \texttt{sat-IN} \\
Sindhi & Devanagari & \texttt{sd} & \texttt{sd-IN} \\
\botrule
\end{tabular*}
\end{table}

\section{User Study Instruments}\label{app:instruments}

\trackb{To be added by the study team.}

\end{appendices}

\bibliography{bibliography-multilingual}

\end{document}
